% 本文件由 LaTeX 排版 Agent 根据有效 translation.json 自动生成。
% author=Basil the Great; work_id=3203; title=论圣神

\chapter{论圣神}

% structure: p0001 source=h1 target=omitted reason=page-title-already-rendered-by-parent

% structure: p0002 source=h2 target=section reason=relative-heading-rank; base=h2

\section{第一章}

\Emph{关于精确考察神学最细微部分之必要性的序言。}\par

1. 我最亲爱的、深受敬仰的兄弟安菲洛奇乌斯，我非常称赞你求知的热忱，以及你勤奋不懈的努力。我极为欣喜地看到，你在表达自己的见解时，表现出如此的谨慎与警醒：在一切言语方式中，论及天主的每一个术语，都不应不作精确考察就放过。你善用了主劝谕中的教诲：\Quote{凡求的，就得到；找的，就寻见。}\BibleRef{路 11:10} 我想，你勤于询问，甚至足以打动最不愿分享的人，使他们分给你他们所拥有的。而这一点更令我敬佩：你没有像我们这个时代大多数人的习气那样，仅仅为了试探而提出问题，而是怀着真诚的愿望去获得实际真理。这个时代不乏挑剔的听众和提问者；但要找到一个渴望求知、以寻求真理作为祛除愚昧之药方的人，是非常困难的。正如猎人的网罗或士兵的埋伏中，诡计通常被巧妙隐藏，大多数提问者的询问也是如此：他们提出论证，并非真要从其中获得什么好处，而是为了在他们未能得到符合自己心意的回答时，似乎就有了争议的合理理由。\par

2. 如果\Quote{愚人求智慧，智慧便算给他，}那么我们当以何等高的价值来估价先知在同一诗句中与\Quote{卓越的谋士}并列的\Quote{智慧的听者}呢？我想，理应以一切赞许来尊重他，并鼓励他进一步前进，分享他的热忱，在他迈向完善的全程中与他一同辛劳。将神学中使用的术语视为首要之事，并努力追踪每一短语、每一音节中的隐藏含义，这是那些在追求真宗教上闲散的人所缺乏的特质，却是所有认识\Quote{我们蒙召的目标}\BibleRef{斐 3:14}的人所具备的；因为我们面前的目标是，在人性所可能达到的程度上，相似于天主。没有知识，就不可能相似；而知识不能没有教训。教训的开端是言语，而音节和词语是言语的部分。因此，考察音节并非离题，也不应因为所提出的问题在某些人看来微不足道而就认为它们不值得关注。真理总是难以追捕的猎物，因此我们必须到处寻找它的踪迹。获取真宗教就像学习手艺一样；两者都是一点一点增长的；学徒不可轻视任何东西。如果一个人轻视最初的基本要素为细小而不屑一顾，他将永远达不到智慧的完善。\par

是和不只是两个音节，然而在这两个小词中常常同时包含着一切美好事物中最好的\Quote{真理}，以及恶所不能逾越的\Quote{谎言}。但何必提是和不是呢？从前，有一位殉道者为基督作证，仅凭点头就被认为已完全履行了真宗教的要求。那么，既然如此，神学中哪个术语如此渺小，以至于它在天平上因使用正确或错误而产生的份量不重要呢？关于法律，我们被告知\Quote{一点一划也绝不会过去}\BibleRef{玛 5:18}；那么，我们忽略甚至最小的细节，怎么能是安全的呢？你自己要求我们彻底筛选的那些点，既小又大。它们的使用只是一瞬间，也许因此被认为微不足道；但是，当我们衡量它们含义的力量时，它们是伟大的。它们可以比作芥菜种，虽然是种子中最小的，但当正确培育并展开其中潜伏的力量时，就长到它固有的高度。\par

如果有人嘲笑我们在音节上的精确（如圣咏所言），让他知道他所收获的是徒劳的果实；而我们，既不屈服于人的指责，也不被他们的轻视打败，继续我们的研究。事实上，我非但不因为这些事微小而感到羞耻，反而即使我能获得其尊贵中的极小一部分，我也要庆幸自己获得了殊荣，并要告诉我的兄弟和共同研究者，他因此获得了不小的收益。\par

因此，虽然我意识到包含在小词中的争论是很大的，但基于对奖赏的盼望，我不畏劳苦，坚信这讨论既对我自己有益，也对我的听众有莫大的好处。所以现在，如果可以这么说的话，在圣神本人的帮助下，我将着手阐述主题；并且，如果你愿意，为了让我进入讨论的正轨，我将暂时回溯我们眼前问题的起源。\par

3. 最近，当我与民众一起祈祷，并使用完整的赞颂词向天主圣父时，我交替使用了两种形式：有时说\Quote{\Emph{偕同}圣子\Emph{及}圣神}，有时说\Quote{\Emph{藉着}圣子\Emph{并在}圣神内}。当时，一些在场的人攻击我，说我在引入新颖且相互矛盾的术语。然而，你主要是为了使他们受益，或者，如果他们完全不可救药，为了那些可能遇到他们之人的安全，表达了应当出版一些关于所用音节背后力量的明确教导的意见。因此，我将尽可能简洁地写作，努力为讨论确立一些公认的原则。\par

% structure: p0010 source=h2 target=section reason=relative-heading-rank; base=h2

\section{第二章}

\Emph{异端者仔细考察音节的起源。}\par

4. 这些人在音节和词语上的吹毛求疵，并非如人们所想的那般简单直接；其所导致的祸害也不小。其中隐藏着一个反对真宗教的深刻而隐蔽的计谋。他们执意争论，声称提及圣父、圣子和圣神的方式是不相同的，仿佛由此就能轻易证明其本性有所差异。他们持有埃提乌斯（\Person{Aetius}）所发明的一个古老诡辩，此人是这一异端的拥护者。在他的某一封信中有一段文字，大意是：本性不同的事物用不同的术语表达；反过来，用不同术语表达的事物在本性上也是不同的。为证明这一论断，他引用了宗徒的话：\BibleQuote{惟一天主圣父，万物都出于祂，……惟一的耶稣基督，万物都藉着祂。} \BibleRef{格前 8:6} 他接着说：\Quote{所以，这些术语彼此之间的关系如何，其所指的本性之间的关系也如何；正如\OriginalTerm{出于祂}{of Whom}与\OriginalTerm{藉着祂}{by Whom}这两个术语不同，圣父与圣子也不同。} 这一异端正是这些人关于所论短语的无聊诡辩所依赖的。因此，他们把\Quote{\OriginalTerm{出于祂}{of Whom}}这句话配给天主圣父，仿佛这是祂特有的一份；他们把\Quote{\OriginalTerm{藉着祂}{by Whom}}这句话限制给天主圣子；把\Quote{\OriginalTerm{在祂内}{in Whom}}这句话归给圣神，并说这些音节的用法绝不互换，为的是——如我已说过的——用语言的变化来指示本性的变化。确实很明显，他们在字眼上的争论实际上是在竭力维护他们不敬神明的论点。\par

他们使用\Quote{\Emph{of} \OriginalTerm{由祂}{whom}}一词，意指创造者；使用\Quote{\Emph{through} \OriginalTerm{藉着祂}{whom}}一词，意指从属的动因或工具；使用\Quote{\Emph{in} \OriginalTerm{在祂内}{whom}}或\Quote{\Emph{in} \OriginalTerm{在其中}{which}}一词，意指时间或地点。所有这些的目的是：宇宙的创造者可能被视为并不高于工具，而圣神可能显得除了来自地点或时间的贡献外，并未给现存事物增添任何东西。\par

% structure: p0014 source=h2 target=section reason=relative-heading-rank; base=h2

\section{第三章}

\Emph{对音节的系统讨论源自外邦哲学。}\par

5. 然而，他们因对外邦作家的深入研究而误入此歧途；这些外邦作家分别将\Quote{\Emph{of} \OriginalTerm{由……}{whom}}和\Quote{\Emph{through} \OriginalTerm{藉着……}{whom}}这两个术语用于本性不同的事物。他们认为，通过\Quote{\OriginalTerm{由……}{of whom}}或\Quote{\OriginalTerm{由……}{of which}}一词，所指示的是质料；而\Quote{\OriginalTerm{藉着……}{through whom}}或\Quote{\OriginalTerm{藉着……}{through which}}一词，则代表工具，或者一般而言，从属的动因。或者更确切地说——因为我们似乎没有理由不详细考察他们的整个论点，并简要揭露其与真理的不相容以及与其自身教导的矛盾——虚妄哲学的研究者在阐述原因的多样性并区分其特殊含义时，将某些原因定义为主要的，某些定义为协同的或共因的，而其他一些则具有\Quote{\OriginalTerm{不可或缺的条件}{sine qua non}}的性质，即必不可少的。\par

对于这些原因中的每一个，他们都有独特而特定的术语用法，以致制造者与工具以不同的方式被指称。对于制造者，他们认为正确的表达是\Quote{\Emph{by} \OriginalTerm{被……}{whom}}，坚称长凳是由木匠\Quote{\OriginalTerm{被}{by}}制造出来的；而对于工具，则是\Quote{\OriginalTerm{藉着……}{through which}}，因为长凳是通过凿子、手钻等器具\Quote{\OriginalTerm{藉着}{through}}制造出来的。同样，他们把\Quote{\OriginalTerm{由……}{of which}}归于质料，因为所造之物是\Quote{\OriginalTerm{由}{of}}木头制成的；而\Quote{\OriginalTerm{按照……}{according to which}}则表示置于工匠面前的设计或样式。因为他要么先在心里打稿，从而将他的构思用于他所要做的事，要么观看事先放在他面前的样式，并据此安排他的工作。他们希望把\Quote{\OriginalTerm{为了……}{on account of which}}这个短语局限于目的或意图，如他们所说，长凳是为了人的使用而制造的。\Quote{\OriginalTerm{在某时/在某处}{in which}}被认为是指时间与地点。它是在何时制造的？在这个时间。在何处？在这个地方。尽管地点和时间对所制造的东西毫无贡献，但没有它们，任何东西都不可能被造出来，因为有效的动因必须同时有地点和时间。正是这些从不切实际的哲学和虚妄的幻觉中得来的细致区分，首先被我们的对手研究并欣赏，然后被转用到有关圣神的简单而纯朴的道理上，以贬低天主圣言，并藐视圣神。甚至连非基督徒作家专门用于无生命工具或最卑微手工服务的短语——我指的是\Quote{\OriginalTerm{藉着或经由……}{through or by means of which}}这个表达——他们也不耻于将其转用于万有之主，而基督徒竟毫无羞愧地将属于锤子和锯子的语言用于宇宙的创造者。\par

% structure: p0018 source=h2 target=section reason=relative-heading-rank; base=h2

\section{第四章}

\Emph{在圣经的用法中，这些音节并无区别。}\par

6. 我们承认，真理之言在许多地方使用了这些表达；然而我们绝对否认圣神的自由受制于异教的狭隘。相反，我们主张圣经根据情况的需要，随场合而变换其表达。例如，\Quote{\Emph{出于}}这个短语并不如他们所假想的那样总是并绝对地指质料，而是更符合圣经的用法，将这一术语应用于至高的原因，正如经上所说：\Quote{一位天主，万物出于祂}\BibleRef{格前 8:6}，又\Quote{万物出于天主}\BibleRef{格前 11:12}。真理之言也曾经常用此术语指质料，如说\Quote{你要用不朽的木造一只方舟}；\Quote{你要用纯金造灯台}\BibleRef{出 25:31}；\Quote{第一个人出于地，属于土}\BibleRef{格前 15:47}；\Quote{你像我一样是泥土所形成}。但是这些人，如我们已指出的，为了建立本性的区别，竟规定这个短语只适用于圣父。这一区分最初源自异教权威，但他们在这里并未表现出忠实的精确限制。对于圣子，他们按照其导师的教导，称之为工具；对于圣神，则称之为场所；因为他们说\Emph{在}圣神内，和\Emph{通过}圣子。但当他们将\Quote{出于}应用于天主时，他们不再追随异教的例子，而是如他们所说，转向宗徒的用法，如经上说：\Quote{但你们出于祂，在基督耶稣内}\BibleRef{格前 1:30}，以及\Quote{万物出于天主}\BibleRef{格前 11:12}。那么，这种系统讨论的结果是什么？有一个原因的属性，一个工具的属性，一个场所的属性。所以圣子在本质上与父不同，如同工具与工匠不同；圣神的区别则如同场所或时间与工具的本性或使用工具者的本性不同。\par

% structure: p0021 source=h2 target=section reason=relative-heading-rank; base=h2

\section{第五章}

\Emph{\Quote{经由谁}也用于圣父，以及\Quote{出于谁}用于圣子和圣神。}\par

7. 在这样描述了对手论证的结果之后，我们现在将继续表明，如我们所提议的，圣父并非先取\Quote{出于谁}然后抛弃\Quote{经由谁}给圣子；并且这些人关于圣子拒绝让圣神分享\Quote{出于谁}或\Quote{经由谁}的裁决是没有真理的，这是他们新奇的措辞分配的限制。\Quote{一位天主，圣父，万物出于祂；一位主耶稣基督，万物经由祂}\BibleRef{格前 8:6}。\par

是的，但这些是一位作者的话，他不是在制定规则，而是在仔细区分\OriginalTerm{位格}{hypostases}。\par

宗徒这样写作的目的不是引入本性的不同，而是表明圣父与圣子的概念不被混淆。这些短语并不彼此对立，也不像战争中排成队列的骑兵那样，使它们所应用的本性相互冲突，这一点从所讨论的段落中非常清楚。蒙福的保罗将两个短语应用于同一个主体，写道：\Quote{万物出于祂，经由祂，并归于祂}\BibleRef{罗 11:36}。即使读者稍微注意一下词义，也会承认这显然指主。宗徒刚刚引用了依撒意亚的预言：\Quote{谁曾知道主的心，或谁曾做过祂的谋士？}然后接着说：\Quote{因为万物出于祂，来自祂，并归于祂。}先知所言指的是天主圣言、万有的创造者，这可以从紧接的前文得知：\Quote{谁曾用手心量过水，用手掌测过天，用斗量过地上的尘土，用秤称过山岭，用天平称过丘陵？谁曾指挥过上主的神，或作祂的谋士教导过祂？}\BibleRef{依 40:12}现在，这段经文中的\Quote{谁}并非指绝对不可能，而是指稀少，正如经文所说：\Quote{谁肯为我起来攻击作恶的人？}以及\Quote{谁人渴望生命？}以及\Quote{谁能上到上主的山？}所讨论的经文也是如此：\Quote{谁曾指挥（七十士译本作'知道'）上主的神，或作祂的谋士知道祂？}\Quote{因为父爱子，把一切事都指给祂看}\BibleRef{若 5:20}。祂是那位持握大地、用手掌握住它，使万物有序并装饰，将山岭安放在其位，测量诸水，赋予宇宙中万物其适当等级，仅用祂的一小部分能力就环抱整个天，先知以比喻称之为手掌。宗徒接着恰当地说：\Quote{万物出于祂，经由祂，并归于祂}\BibleRef{罗 11:38}。因为万有存在的原因出于祂，按照天主圣父的旨意。万有经由祂得以持守和构成，因为祂创造了万有，并给每一事物分配维持其健康与保存所需。因此万有转向祂，带着不可抗拒的渴望和难以言喻的爱恋，注视着他们生命的\Quote{创始者}\BibleRef{宗 3:15}和维持者，正如经上所写：\Quote{众人的眼睛都是向您看}，又：\Quote{这一切都向您看}，\Quote{您张开手，满足一切生物的愿望}。\par

8. 但如果我们的对手反对我们这种解释，有什么论据能救他们不陷入自己的陷阱呢？\par

\Quote{从他}、\Quote{藉着他}、\Quote{归向他}这三个表述是用在主身上，那么它们就只能用于天主圣父。于是他们的规则无疑会失败，因为我们不仅发现\Quote{从他}也用于圣父，而且\Quote{藉着他}也用于圣父。如果后一个短语并不表示任何贬义，那为什么它偏偏要被限制用于圣子，仿佛传达了从属的含义？如果它总是且处处暗示服务，那么让他们告诉我们，荣耀的天主和基督的父是从属于哪个上级。\par

他们就这样被自己推翻，而我们的立场则从两方面得到确认。假设经文被证明是指圣子，那么\Quote{从他}将被发现适用于圣子。另一方面，如果坚持先知的言论与天主有关，那么就会承认\Quote{藉着他}也合宜地用于天主，并且两个短语具有同等价值，因为两者都以同等的效力用于天主。在任一假设下，这两个术语既然用于同一位格，便将被证明是等价的。但让我们回到我们的主题。\par

9. 在致厄弗所人的书信中，宗徒说：\Quote{惟用爱德持守真理，在各方面长进，归于他，他是元首，就是基督；全身都靠他联络得合式，百节各按各职，照着各肢体的功用，彼此相助，便叫身体渐渐增长。} \BibleRef{弗 4:15}\par

再者在致哥罗森人的书信中，对于那些不认识独生者的人，提到了那持守\Quote{元首}的，即基督，\Quote{全身靠关节和韧带得滋养供应，以天主的增长而增长。} \BibleRef{哥 2:19} 而且我们在另一段经文中得知基督是教会的头，宗徒说\Quote{使他为教会作万有之首}，\BibleRef{弗 1:22} 并说\Quote{从他的满盈中我们都领受了}，\BibleRef{若 1:16}。主自己说\Quote{他要从我领受，并要传告给你们}，\BibleRef{若一 16:15}。总之，细心的读者会察觉到\Quote{从他}以多种方式使用。例如，主说\Quote{我感觉到有能力从我出去}，同样我们经常看到\Quote{从他}用于圣神。\Quote{那为圣神撒种的}，经上说，\Quote{将从圣神收获永生}，\BibleRef{迦 6:8}。若望也写道：\Quote{我们借此知道他在我们内，是由\OriginalTerm{由}{\GreekText{ἐκ}}他赐给我们的圣神}，\BibleRef{若一 3:24}。\Quote{她所怀的}，天使说，\Quote{是由圣神而来的}，\BibleRef{玛 1:20}，主说\Quote{由圣神而生的就是神}，\BibleRef{若 3:6}。以上就是目前的情况。\par

10. 现在必须指出，\Quote{藉着他}这一短语被圣经同等用于圣父、圣子和圣神。证明用于圣子固然繁琐，不仅因为这事众所周知，更因为我们的对手也强调这一点。我们现在证明\Quote{藉着他}也用于圣父。经上说：\Quote{天主是信实的，你们是凭\OriginalTerm{凭他}{\GreekText{δι᾿} \GreekText{οὖ}}蒙召，与他圣子有分}，\BibleRef{格前 1:9}；以及\Quote{保禄因天主的旨意\OriginalTerm{因}{\GreekText{διά}}作耶稣基督的宗徒}；又说：\Quote{所以你不再是奴仆，而是儿子；如果是儿子，也藉着天主成为继承人。} 以及\Quote{正如基督因圣父的荣耀\OriginalTerm{因}{\GreekText{διά}}从死者中复活。} 此外，依撒意亚说：\Quote{祸哉，那些深藏谋划而不藉着主的人们。} 关于圣神使用这一短语的许多证据也可以引用。经上说：\Quote{天主藉着圣神\OriginalTerm{藉着}{\GreekText{διά}}将他启示给我们}，\BibleRef{格前 2:10}；另一处说：\Quote{你要藉着圣神\OriginalTerm{藉着}{\GreekText{διά}}保守所托付你的善事}，\BibleRef{弟后 1:14}；又说：\Quote{这人藉着圣神\OriginalTerm{藉着}{\GreekText{διά}}得了智慧的言语}，\BibleRef{格前 12:8}。\par

11. 同样，也可以说\Quote{在……内}这个词，圣经承认它用于天主圣父。在旧约中，经上说：\Quote{我们藉着\OriginalTerm{藉着}{\GreekText{ἐν}}天主奋力行事}；以及\Quote{我的赞美将不断地在\OriginalTerm{在}{\GreekText{ἐν}}你内}；又说\Quote{我要因你的名\OriginalTerm{因}{\GreekText{ἐν}}欢欣}。在保禄书信中，我们读到：\Quote{天主在\OriginalTerm{在}{\GreekText{ἐν}}内创造万有}，\BibleRef{弗 3:9}；以及\Quote{保禄、息耳瓦诺和弟茂德，致书给德撒洛尼人在天主父内的\OriginalTerm{在}{\GreekText{ἐν}}教会}，\BibleRef{得后 1:1}；以及\Quote{也许我因天主的旨意\OriginalTerm{因}{\GreekText{ἐν}}终能顺利去到你们那里}，\BibleRef{罗 1:10}；以及\Quote{你以天主\OriginalTerm{以}{\GreekText{ἐν}}夸口}，\BibleRef{罗 2:17}。实例实在数不胜数；但我们所需要的与其说是展示大量证据，不如说是证明我们对手的结论是站不住脚的。因此，我将省略证明这一用法用于我们的主和圣神，因为这是众所周知的。但我不能不指出，\Quote{明智的听者}将通过对立的方法找到当前命题的充分证据。因为如果语言的差异，正如我们被告知的，表明本性已经改变，那么就让语言的一致迫使我们的对手羞愧地承认本质没有改变。\par

12. 不仅是在神学中用语的使用有所变化，而且每当一个词取另一个词的意义时，我们常常发现它们从一个主体转移到另一个主体。例如，亚当说：\Quote{我\OriginalTerm{\Emph{藉着}}{through}天主获得了一个人}，其含义与\Quote{\OriginalTerm{\Emph{从}}{from}天主}相同；在另一处经文说：\Quote{梅瑟奉主的命令……吩咐以色列}，又说：\Quote{解释岂非出自天主？}若瑟向囚犯讲解梦境时，不说\Quote{\Emph{从}天主}，而直接说\Quote{\Emph{藉着}天主}。反过来，保禄在使用\Quote{\Emph{从}他}代替\Quote{\Emph{藉着}他}时，说：\BibleQuote{生于女人}\BibleRef{迦 4:4}。他在另一处明确区分了这两种用法，指出女人由男人而出，男人则藉着女人而生，正如所说：\BibleQuote{因为女人是由男人而出，男人也是藉着女人而生}\BibleRef{格前 11:12}。然而在讨论的经文中，宗徒在说明用法的多样性时，\Emph{顺便}纠正了那些认为主的身体是属灵身体之人的错误，并为了显示承载天主的血肉是由普通人性团块形成，而强调了更明确的介词。\par

\Quote{\OriginalTerm{藉着女人}{through a woman}}这一说法可能会引起对生育中纯粹传递的怀疑，而\Quote{\OriginalTerm{从女人}{of a woman}}则令人满意地表明了母亲与后代共享本性。宗徒绝非自相矛盾，而是表明了这些词语可以毫无困难地互换。因此，既然\Quote{\OriginalTerm{从他}{from whom}}这一词语被转移到了那些被判定为正确使用\Quote{\OriginalTerm{藉着他}{through whom}}的同一主体上，那么将这些短语一成不变地加以区分，以便对真宗教进行错误指责，这又有什么一致性可言呢？\par

% structure: p0035 source=h2 target=section reason=relative-heading-rank; base=h2

\section{第六章}

\Emph{与那些主张子不与父同在，而在父之后之人的论战。关于同等光荣。}\par

13. 我们的对手如此巧妙而乖戾地对抗我们的论证，甚至无法以无知为借口。显然他们对我们因将\Quote{光荣颂}与父一起归于独生子，以及不将圣神与子分开而感到恼怒。为此，他们称我们为革新者、革命者、新词创造者，以及其他一切可能的侮辱称号。但我远未被他们的辱骂所激怒，若非他们的失落使我感到\Quote{忧愁和不断的痛苦}，我几乎可以说我感谢他们的亵渎，似乎他们是为我提供祝福的媒介。因为经上说：\BibleQuote{几时人为了我的缘故辱骂你们，你们是有福的}\BibleRef{玛 5:11}。他们愤慨的理由如下：根据他们的说法，子不与父同在，而在父之后。因此，光荣应\Quote{\OriginalTerm{\Emph{藉着他}}{through him}}归于父，而不是\Quote{\OriginalTerm{\Emph{与他一起}}{with him}}归于父；因为\Quote{与他一起}表达尊严的平等，而\Quote{藉着他}表示从属关系。他们还断言，圣神不应与父和子并列，而应处于子和父之下；不应并列，而应从属；不应同列，而应次列。\par

他们用这类技术术语歪曲了信仰的单纯和朴实，并藉其机巧，使他人无法保持无知，从而为自己切断了无知可能要求的借口。\par

14. 让我们首先问他们这个问题：他们说子\Quote{在父之后}是什么意思？是在时间之后、顺序之后，还是尊严之后？但在时间上，没有人会愚蠢到声称万代的创造者居于第二位，因为在父与子的自然结合中没有间隔。而且，就我们对人类关系的理解而言，不可能认为子晚于父，这不仅因为父与子根据他们之间的相互关系被共同设想，而且因为时间上的在后是指离现在较近的主体，而先前是指较远的主体。例如，诺厄时代所发生的事先于索多玛人所发生的事，因为诺厄离我们更远；而索多玛人的事是后来发生的，因为它们似乎更接近我们的时代。但是，除了违反真宗教之外，用距现在的远近去衡量超越一切时间和万代的生命的存在，难道不是极端的愚蠢吗？就好像天主父能与存在于万代之前的天主子相比并超越他，正如那些有起始和朽坏的事物被描述为彼此先后一样？\par

父的超越遥远实际上是不可想象的，因为思想和理智完全无法超越主的受生；\Person{圣若望}用两个词巧妙地将这一概念限制在界限之内：\BibleQuote{在\Emph{起初已有}圣言}。因为思想无法超越\Quote{已有}，想象也无法超越\Quote{\Emph{起初}}。让你的思想无论怎样回溯，你都无法超越\Quote{已有}；无论你如何努力想看到子之外的东西，你都会发现不可能超越\Quote{\Emph{起初}}。因此，真宗教教导我们如此思想子与父同在。\par

15. 若他们确实认为圣子相对于圣父有一种降格，好像祂处于较低的位置，以至于圣父坐在高处，而圣子被推到下面的次座，就让他们承认他们的意思吧。我们不再多说。一个明确的陈述将立即暴露其荒谬。那些拒绝承认圣父充满万有的人，甚至不能维持他们论证的逻辑连贯性。健全者的信德是：天主充满万有；\BibleRef{弗 4:10}但那些在圣父和圣子之间划分上下的人，甚至不记得先知的话：\Quote{我若升到天上，你在那里；我若下到阴间，你也在那里。}现在，暂且忽略所有那些将位置归于无形之物的无知证据，他们有何借口攻击圣经，其敌对如此无耻，正如经文\Quote{你坐在我右边}和\Quote{坐在天主威能的右边}？\Quote{右边}一词并非如他们争辩的那样指示较低的位置，而是表示关系的平等；它不能被按肉体理解，因为那样天主会有左边，但圣经用表示尊荣的庄严语言来呈示圣子尊严的宏伟。因此，我们的对手只能声称这个表述表示等级的低下。让他们学习\Quote{基督是天主的德能和天主的智慧}\BibleRef{格前 1:24}，以及\Quote{他是不可见的天主的肖像}\BibleRef{哥 1:15}，和\Quote{是他光荣的辉耀}\BibleRef{希 1:3}，并且\Quote{天主圣父印证了他}\BibleRef{若 6:27}，将祂自己烙印在祂身上。\par

现在我们该称这些段落以及类似它们贯穿整部圣经的经文，是屈辱的证明，还是独生子的威严及其与圣父光荣平等的公开宣告？我们请求他们听从主自己明确陈述祂与圣父光荣的同等尊严的话：\Quote{看见了我就是看见了父；}\BibleRef{若 14:9}以及\Quote{当圣子在他父的光荣中来临时}\BibleRef{谷 8:38}；\Quote{应尊敬圣子如同尊敬圣父}\BibleRef{若 5:23}；\Quote{我们见过他的光荣，正如父独生子的光荣}\BibleRef{若 1:14}，以及\Quote{在父怀里的独生天主。}他们不顾所有这些段落，却将圣子置于为祂仇敌预备的地方。父的怀里是适合圣子的座位，但脚凳之处是为那些必须被降伏的人。\par

我们只是粗略地触及这些证据，因为我们的目标是转到其他要点。您可以在闲暇时把这些证据收集起来，然后默观独生子光荣的高度和德能的卓越。然而，对于善意听者而言，这些也并非不重要，除非\Quote{右边}和\Quote{怀里}这些词被按肉体和贬义来接受，从而立即将天主局限于地方界限内，并发明形式、形状和身体姿势，所有这些都与绝对、无限、无形体概念完全无关。此外，有一个事实：其中贬义的概念对圣父和圣子而言是相同的；因此，无论谁重复这些论点，都不会剥夺圣子的尊严，却会招致亵渎圣父的指控；因为人对圣子所犯的任何妄用，都不能不转移到圣父身上。如果他以优先的方式将上位分配给圣父，并断言独生圣子坐在下面，他将发现他想象的产物附带着身体的所有相应条件。如果这些是醉酒妄想和疯狂谵妄的产物，那么对于被主亲自教导\Quote{不尊敬圣子的，就是不尊敬圣父}\BibleRef{若 5:23}的人而言，拒绝敬拜和光荣那位在本性、光荣和尊严上与圣父结合的圣子，岂能符合真宗教？我们将说什么？在可怕的普世万有审判之日，我们将有什么正当的辩护？如果主清楚宣告祂将\Quote{在他父的光荣中}\BibleRef{玛 16:27}来临，如果斯德望看见耶稣站在天主的右边\BibleRef{宗 7:55}，如果保禄在圣神内为基督作证\Quote{他在天主的右边}\BibleRef{罗 8:34}，如果圣父说\Quote{你坐在我右边}，如果圣神见证他已经坐在\Quote{天主权能的右边}\BibleRef{希 8:1}，我们却试图将这同享光荣和宝座者从平等的状态贬低到较低的状态？我认为，站和坐表示本性的固定和完全稳定，正如巴路克想要显示天主的存有方式的不变和不动时所说：\Quote{因为你永远坐着，而我们永远消亡。}此外，右边的地方在我看来表示光荣的平等。因此，剥夺圣子参与光荣颂是鲁莽的，好像他只配被排在较低的光荣地位。\par

% structure: p0044 source=h2 target=section reason=relative-heading-rank; base=h2

\section{第七章}

\Emph{反对那些断言对圣子说\Quote{偕同祂}是不合适，而正确的说法是\Quote{借着祂}的人。}\par

16. 但他们争辩说，使用\Quote{\OriginalTerm{与祂}{with him}}这个说法是完全陌生和不寻常的，而\Quote{\OriginalTerm{藉着祂}{through him}}在圣经中最为常见，在弟兄们的语言中也很普遍。我们对此如何回答？我们说：没有听过你们的耳朵有福，没有被你们言语的伤害所保全的心有福。对你们这些爱基督的人，我却说，教会承认两种用法，而不排斥任何一种，以为它破坏另一种。因为每当我们默想独生子的本性的尊威和祂位格的高贵时，我们就见证荣耀是\Emph{与}圣父同在的；而每当我们想到祂赐予我们的恩赐，以及我们进入天主的家室，被接纳为家人时，我们就承认这恩宠是\Emph{藉}着祂并\Emph{因}着祂为我们实现的。\par

由此可见，在荣耀颂中适合使用\Quote{\OriginalTerm{与祂}{with whom}}这个说法，而在感恩时则特别适合使用\Quote{\Emph{藉}着祂}这个说法。而断言\Quote{与祂}这一说法在虔诚者的用法中不熟悉，也是完全不真实的。凡是性格健全，认为古代的尊贵比新奇的新颖更可敬，并保存了他们祖传的传统纯正无伪的人，无论在城市还是乡村，都使用这个说法。正相反，那些厌弃熟悉和惯常之事，并傲慢地攻击古老为陈腐的人，却欢迎创新，正如在衣着上，喜好炫耀的人总是偏爱某种全新之物而非平常穿着的。所以你甚至今天仍可看到乡间人的语言保留着古老的风格，而我们这些精于\OriginalTerm{词语之争}{logomachy}的狡猾专家，其语言却打着新哲学的烙印。\par

我们父辈所说的，我们同样说：圣父和圣子的荣耀是共同的；因此我们以\Quote{与圣子在一起}的方式向圣父献上光荣颂。但我们不只依靠这是父辈的传统；因为他们也是遵循圣经的含意，并从我在前面几句话中从圣经引出并摆在你们面前的证据出发。因为\Quote{光辉}总是与\Quote{荣耀}一同被思想，\Quote{肖像}与原型一同被思想，\BibleRef{格后 4:4}，圣子总是且处处与圣父在一起；连名字的紧密联系，更不用说事物的本性，都不容分离。\par

% structure: p0049 source=h2 target=section reason=relative-heading-rank; base=h2

\section{第八章}

\Emph{\Quote{藉着谁}有几种用法；以及\Quote{与谁}在何种意义上更合适。说明圣子如何接受命令，以及祂如何被派遣。}\par

17. 因此，当宗徒\BibleQuote{藉着耶稣基督感谢天主}，\BibleRef{罗 1:8}，并且又说\BibleQuote{藉着祂}我们\BibleQuote{领受了恩宠和宗徒的职务，为使万邦服从信仰}，\BibleRef{罗 1:5}，或\BibleQuote{藉着祂我们获得了进入这恩宠的门路，在这恩宠中我们站立并且喜乐}，\BibleRef{罗 5:2}，他陈述了圣子赐予我们的恩惠，有时使美善恩赐的恩宠从圣父传递给我们，有时又通过祂自己把我们带到圣父面前。因为他说\BibleQuote{藉着祂我们领受了恩宠和宗徒的职务}，\BibleRef{罗 1:5}，他宣布美善恩赐的供应源自那源头；又说\BibleQuote{藉着祂我们获得了进入的门路}，\BibleRef{罗 5:2}，他陈述了我们的接纳并藉着基督成为\Quote{天主的家室}。那么，承认祂为我们所施的恩宠是否减损了祂的荣耀？难道不是说述说祂的恩惠是光荣祂的恰当理由吗？正是为此，我们发现圣经不是用一个名字描述主给我们，甚至也不是只用那些表明祂天主性和尊威的术语。有时它使用描述祂本性的术语，因为它承认\BibleQuote{那超越一切名号的名字}，\BibleRef{斐 2:9}，圣子的名号，并提到真圣子，独生天主，天主的德能，和智慧，\BibleRef{格前 1:24}，以及圣言。然后，因为恩赐给我们的恩宠有多种方式，祂因良善的富饶，按照祂丰盛的智慧，\BibleRef{弗 3:10}，施予需要的人，圣经用无数其他称号称呼祂，称祂为牧人、君王、医师、新郎、道路、门、泉源、食粮、斧头和磐石。这些称号并不表明祂的本性，而如我所评论的，表明祂出于对自己受造物的温柔之心，按照各人的特别需要，施予需要者的各种有效运作。那些逃到祂治理的照顾下，并通过坚忍悔改了自己任性行为的人，祂称为\Quote{羊}，并承认自己对于听祂声音、拒绝理会异端教导的人，是\Quote{牧人}。因为祂说：\BibleQuote{我的羊听我的声音。}对于已经达到更高阶段、需要正义王权的人，祂是君王。由于通过祂诫命的直路，祂引导人行善，又因为祂安全地关闭所有因信德投奔祂以寻求更高智慧祝福的人，祂是门。\par

因此祂说：\BibleQuote{凡由我进入的，必得安全，可以自由出入，并能找到牧场。}\BibleRef{若 10:9} 再者，因为对信者而言，祂是坚固、不动摇、比任何壁垒更难摧毁的防御，祂是磐石。在这些称号中，当祂被称为门或道路时，\Quote{藉着祂}这个短语非常合适且明白。然而，作为天主与圣子，祂与圣父同受光荣，因为\BibleQuote{因耶稣之名，一切在天上的、地上的和地底下的，无不屈膝，并且一切唇舌无不宣认耶稣基督是主，为光荣天主圣父。}\BibleRef{斐 2:10} 因此我们同时使用两个表述，一个表达祂本身的尊严，另一个表达祂对我们的恩宠。\par

18. 因为\Quote{藉着祂}获得我们灵魂的一切援助，并且按照每一种照顾，已经想出了一个适当的称号。所以当祂将无可指摘、没有瑕疵或皱纹的灵魂\BibleRef{弗 5:29}如同纯洁的少女献给自己时，祂被称为新郎；但每当祂从魔鬼的邪恶打击中接收一个遭受严重痛苦的人，在其罪的沉重软弱中医治他时，祂被称为医师。难道祂对我们的这种照顾竟会贬低我们对祂的思念吗？或者说，它反而会激起我们对救主的大能和仁爱的惊叹，因为祂既忍受了与我们一同受苦于我们的软弱，又能屈尊就卑于我们的脆弱？因为不是天、地、大海，不是水中和陆地上的生物，不是植物、星辰、空气、季节，也不是宇宙秩序中的巨大多样性，如此清楚地表明祂权能的卓越，如同那不可理解的天主，竟然能无痛苦地藉着肉身与死亡激烈交锋，以便藉着祂自己的苦难赐予我们免受痛苦的恩赐。的确，宗徒说：\BibleQuote{藉着那爱我们的主，我们在这一切事上大获全胜。}\BibleRef{罗 8:37} 但是在这种表述中，没有任何卑下和从属服役的暗示，而是藉着\BibleQuote{祂大能的威力}所提供的援助\BibleRef{弗 6:10}。因为祂亲自捆绑了壮士，夺取了他的财物，即我们这些人类，我们的仇敌曾在每一种邪恶活动中滥用我们，并使我们成为\BibleQuote{合乎主用的器皿}\BibleRef{弟后 2:21}，我们通过预备我们自身可控制的部分而为每一项工作得以成全。这样，我们便藉着祂得以亲近圣父，从\BibleQuote{黑暗的权势}被转移到\BibleQuote{光明中圣徒的基业}中\BibleRef{哥 1:12}。然而，我们不应该把藉着圣子的\OriginalTerm{救恩计划}{œconomy}看作是因奴隶的卑贱地位而被迫的从属服役，而是出于善良和怜悯，按天主圣父的旨意，为祂自己的受造物有效地工作的自愿关怀。因为如果我们承认在一切由祂在过去和现在所完成的事中，都显示出祂权能的圆满，并且绝不将之与圣父的旨意分离，那么我们就与真宗教一致。例如，每当主被称为道路时，我们就被引向更高的意义，而不是来自该词世俗含义的意义。我们通过道路理解为那种逐步、有秩序地经由正义的作为和\Quote{知识的照耀}而朝向成全的进展；不断渴求前面的事，伸向那些仍存的事，直到我们达到蒙福的终点，即认识天主——主藉着祂自己将这认识赐予那些信赖祂的人。因为我们的主是一条本质上是善的道路，不知迷失与歧途，通向本质上是善的圣父。因为祂说：\BibleQuote{除非藉着我，没有人能到父那里去。}\BibleRef{若 14:6} 这就是我们\Quote{藉着圣子}上达于天主的道路。\par

19. 接下来，我们应当按顺序指出圣父\Quote{藉着祂}所赐予我们的恩惠之品的特征。既然一切受造的本性，无论这个可见世界还是头脑中所想到的一切，若没有天主的眷顾和上智安排就无法持存，那么创造者圣言、\OriginalTerm{独生天主}{Only begotten God}，按着各者需要之度量分配祂的救助，由于祂恩赐的领受者种类和品格繁多而分施各种多样的仁慈，却又合乎个体需求之必需。那些被囚于愚昧黑暗中者，祂光照他们：为此祂是真光。\BibleRef{若 1:9} 按行为应得的报应分施回报，祂审判：为此祂是公义的审判者。\BibleRef{弟后 4:8} \Quote{因为圣父不审判任何人，但把一切审判都交给了圣子。}\BibleRef{若 5:22} 那些从生命的高峰堕落于罪者，祂从他们的失足中扶起他们：为此祂是复活。\BibleRef{若 11:25} 以祂权能的触摸和祂良善的旨意有效地运作，祂成就一切。祂牧放；祂光照；祂滋养；祂医治；祂引导；祂举起；祂呼唤那不存在的成为存在；祂维系那已被创造的。如此，来自天主的善事\Quote{藉着圣子}到达我们，祂在每件事上的运作比言语所能述说的更快。因为闪电不及，空中光行不及，眼目锐转不及，我们思绪的移动也不及。不，藉着神圣的动能，这些中的每一个在速度上都被远远超越，如同一切活物中最慢的在运动中不及飞鸟，甚或风，或天体的飞驰：或者，且不提这些，连我们自己的思绪也如此。因为那位\Quote{以自己大能的圣言支撑万有}\BibleRef{希 1:3}且不藉身体工具运作，也不需手助以形成和塑造，却以顺服的追随和不强迫的掌管万有之本性者，需要何等的时间呢？正如\Person{友弟德}所说：\Quote{祢思想过，凡祢决定的事都已齐备。}另一方面，免得我们因所成就之工的伟大而被牵引去想象主是无始的，那\OriginalTerm{自有者}{Self-Existent}说了什么？\Quote{我藉着父而生活[by, A.V.]，}\BibleRef{若 6:57}以及天主的能力；\Quote{圣子自己什么也不能做[can, A.V.]。}\BibleRef{若 5:19}而那\OriginalTerm{完备的智慧}{self-complete Wisdom}呢？我领受了\Quote{一条命令，该说什么，该讲什么。}\BibleRef{若 12:49}通过这一切话，祂正引导我们认识圣父，并将我们对一切受造之物的惊奇归向祂，为使我们\Quote{藉着祂}能认识圣父。因为圣父不是从运作的差异而被认识，不是藉着显示一种分离而特殊的能力；因为无论祂看见圣父做什么，\Quote{圣子也照样做；}\BibleRef{若 5:19}但祂享受我们对一切因从\OriginalTerm{独生者}{Only Begotten}而来的荣耀而发生的惊奇，既因行这事者本身也因事工的伟大而喜悦，并被所有承认祂为我们主耶稣基督之圣父的人高举，\Quote{万有是藉着祂[by whom, A.V.]，万有也是为了祂。}因此，主说：\Quote{我的一切都是祢的，}\BibleRef{若 17:10}仿佛受造之物的主权被赐予祂；\Quote{祢的一切都是我的，}仿佛创造之源从那里来到祂。我们不应设想祂在行动中使用了帮助，或受托以详细委托去执行每一项个别工作的职分——这是一种明显卑贱且完全不符神圣尊贵的状态。相反，圣言充满祂圣父的美善；祂从圣父发出，并依照生祂者的样式行一切。因为若在本质上祂是不变的，那么在能力上祂也是不变的。而那些能力相等者，其运作也全然相等。\Quote{并且基督是天主的能力、天主的智慧。}\BibleRef{格前 1:24} 因此\Quote{万有都是藉着祂[by, A.V.]而造的，}\BibleRef{若 1:3} 以及\Quote{万有都是藉着祂[by, A.V.]并为了祂而受造的，}\BibleRef{哥 1:16} 并非履行任何奴隶性的服务，而是为完成圣父作为创造者的旨意。\par

20. 当祂说：\Quote{我没有凭我自己说话，}\BibleRef{若 12:49} 又说：\Quote{父怎样告诉了我，我就照样讲论，}\BibleRef{若 12:50} 以及\Quote{你们所听到的话，不是我的，而是派遣我来的父的，}\BibleRef{若 14:24} 并在另一处说：\Quote{父怎样命令了我，我就照样作，}\BibleRef{若 14:31} 这不是因为祂缺乏自主的目的或发起的权能，也不是因为祂必须等候预先约定的基调，才使用这种言语。祂的目的是要表明祂自己的意志与父的意志有着不可分割的结合。因此，我们不要理解所谓\Quote{命令}是由发声器官传达的强制性指令，如同对属下一样，吩咐子祂应当做什么。反而要以适合天主性的方式，领悟一种意志的传递，如同物体在镜子中的反射，不受时间影响地从父传到子。因为\Quote{父爱子，并把一切显示给祂，}\BibleRef{若 5:20} 以致\Quote{父所有的一切}都属于子，不是渐渐地一点一点地加给祂，而是与祂一起全部且立刻地拥有。在人当中，一个受过充分技艺训练、并经过长期培训而具有确定且牢固经验的工匠，能够依照他现在所储备的科学方法，独自为将来工作。难道我们竟以为创造万有的天主的智慧，祂是永远完美的、有智慧的、无需导师的、天主的德能，\Quote{在祂内隐藏着一切智慧和知识的宝藏，}需要零碎地教导来划定祂作为的方式和尺度吗？我猜想，在你计算中的虚妄，你打算开办一所学校；你要让一位坐在教师的位子上，另一位则站在门徒的无知中，通过一点一点地给予祂的教训，逐渐学习智慧并迈向成全。因此，如果你有理智遵循逻辑的结果，你会发现子永远受教，却永不能达到成全的终点，因为父的智慧是无限的，而无限的终点是无法把握的。其结果就是，凡不肯承认子从起初就拥有一切的人，绝不会承认祂将达到成全。但我为因论证过程而沦落到的这种卑下观念感到羞愧。因此，让我们回到我们讨论的更崇高的主题吧。\par

21. \Quote{看见我的，就是看见了父；}\BibleRef{若 14:9} 不是精确的肖像，也不是形状，因为天主性不容许组合；而是意志的良善，它与本质同时存在，在父和子中被视为相似、相等，甚至同一。\par

那么，所谓\Quote{成了服从的}是什么意思？\BibleRef{斐 2:8} 所谓\Quote{为我们众人把祂交出来}是什么意思？\BibleRef{罗 8:32} 意思是子从父那里获得祂为了人类而施行良善的能力。但你们还必须听这些话：\Quote{基督救赎了我们脱离了法律的咒诅；}\BibleRef{迦 3:13} 以及\Quote{当我们还是罪人的时候，基督就为我们死了。}\BibleRef{罗 5:8}\par

也要十分注意主的话，留意每当祂向我们教导有关祂的父时，祂习惯使用具有个人权威的措辞，说：\Quote{我愿意，你洁净了吧；}\BibleRef{玛 8:3} 和\Quote{安静，平静了吧；}\BibleRef{谷 4:39} 以及\Quote{但我告诉你们；}和\Quote{你这又聋又哑的鬼，我命令你；}\BibleRef{谷 9:25} 以及其他所有同类表达，为使我们借此认出我们的师傅和创造者，并借前者认识我们师傅和创造者的父。因此，从各方面都证明这真确的教导：父通过子创造，这既不是使父的创造不完美，也不是显示子的行动力软弱，而是表明意志的合一；因此\Quote{借着祂}这一表达包含了对先在原因的承认，而不是用来反对动力因。\par

% structure: p0059 source=h2 target=section reason=relative-heading-rank; base=h2

\section{第九章}

\Emph{符合圣经教导的关于圣神的确定概念}\par

22. 现在让我们探究我们关于圣神的共同概念，包括我们从圣经中收集的关于祂的概念，以及我们从教父们未成文的传统中领受的概念。首先，我们问：谁在听到圣神的称号时，心灵不提升，谁不将自己的思想提升到至高的本性？祂被称为\Quote{天主之神}、\BibleQuote{从父发出的真理之神}\BibleRef{若 15:26}、\Quote{正直之神}、\Quote{引领之神}。祂专属的称号是\Quote{圣神}；这个名称特别适用于一切无形体、纯粹非物质、不可分割的事物。因此，当我们的主教导那位以为天主是地方性崇拜对象的女人的时候，说无形体是不可理解的，祂说\BibleQuote{天主是神}。\BibleRef{若 4:24}因此，当我们听到\Quote{神}时，不可能形成一种受限制、易变化、或被造物般的本性观念。我们被迫将思想提升到至高处，思想一种理智的本质，能力无限，量度无限，不受时间或世代度量，慷慨赐予美善恩赐，一切需要圣化的事物都转向祂，一切在美德中生活的事物都追求祂，仿佛被祂的呼吸所滋润，被帮助走向自然和正确的终向；成全一切其他事物，但祂自身毫无欠缺；生活不是需要恢复，而是生命的供给者；不因增加而增长；而是立刻充满、自立、无处不在、圣化的源头、理智可见的光，仿佛通过祂自身，为每一官能提供光照，以寻求真理；本性不可接近，因良善而被理解，以祂的能力充满万有，但只传给配得的人；不是按一种度量分享，而是按\BibleQuote{信德的比例}\BibleRef{罗 12:6}分配祂的德能；本质单纯，德能多样，完全临在于每一事物，又完全无处不在；无痛苦地分开，分享而不失完整，类似于阳光，其恩泽之光落在享受它的人身上，仿佛只为他照耀，却又照亮大地和海洋，与空气混合。同样，圣神对于每一位接受祂的人，仿佛只赐给他人，却又为全人类发出充分而圆满的恩宠，并被所有分享祂的人享受，按他们的能力，不是按祂的能力，而是按他们的本性。\par

23. 现在，圣神不是通过地方性的相近而与灵魂建立亲密关系。怎么可能有形体接近无形体呢？这种关系来自于克服情欲，这些情欲是后来逐渐因灵魂与肉体的友谊而进入灵魂，使灵魂疏远了与天主的亲密关系。只有当人从因邪恶而沾染的污秽中被净化，恢复到自然之美，如同擦净皇家肖像，恢复其原有形式，只有这样，他才有可能接近护慰者。而祂，如同太阳，借助你洁净的眼睛，在祂自身中向你显示不可见者的肖像，并在这肖像的有福景象中，你将看到原型那不可言喻的美。通过祂的帮助，心灵被提升，软弱者被扶持，前进者得以完善。照耀在从一切污点中洁净的人身上，祂使他们通过与祂的共融而成为属神的。正如当阳光落在明亮透明的物体上时，它们自身也变得光辉，并从自身发出新的光辉，同样，圣神所居住的灵魂，被圣神光照，自身也成为属神的，并将恩宠施予他人。由此，预知未来，理解奥秘，领悟隐藏之事，分配美善恩赐，天上公民身份，加入天使歌咏团，无尽的喜乐，常在天主内，相似天主，以及最高者——成为神。这些，只是众多概念中的少数几个，是从圣神自己的神谕中学到的关于圣神的伟大、尊严和行动的教义。\par

% structure: p0063 source=h2 target=section reason=relative-heading-rank; base=h2

\section{第十章}

\Emph{反对那些说圣神不应与圣父和圣子同等同列的人。}\par

24. 然而，我们必须前进攻击我们的对手，试图驳斥那些反对我们的\Quote{反对论点}，这些论点是从自称的\Quote{伪知识}中衍生出来的。\par

他们断言，圣神不应与圣父和圣子同等同列，因为祂的本性不同，尊严较低。对此，应以宗徒的话来回答：\BibleQuote{我们应当服从天主，而不是人}\BibleRef{宗 5:29}\par

因为我们的主在命令救恩的圣洗时，嘱托祂的门徒\BibleQuote{因父、及子、及圣神之名}\BibleRef{玛 28:19}给万民授洗，并未不屑于与圣神共融，而这些人却声称我们不应将祂与圣父和圣子同等同列，难道不是很明显他们公然违背天主的命令吗？如果他们否认这种并列是表明任何共融与结合，让他们告诉我们为什么我们应该持这种观点，以及他们有什么更紧密的结合方式。\par

如果主确实没有在圣洗中将圣神与圣父和祂自己结合，那么就不要将结合的责任归咎于我们，因为我们既不持有也不说出任何不同的话。相反，如果在圣洗中圣神与圣父和圣子结合，并且没有人如此无耻地说别的话，那么就不要责备我们跟随圣经的话。\par

25. 但所有战争的装备都已对准我们；每一件智力的投射物都已瞄准我们；如今亵渎者的舌头射中又射中，比当年杀害基督的人击打\Person{斯德望}更猛烈。不要让他们成功隐瞒这一事实：虽然攻击我们只是战争的借口，这些行动的真实目标更高。他们说，他们正在准备器械和陷阱来攻击我们；他们正互相呼喊，根据各人的力气或狡诈，要上前来。但攻击的目标是信德。所有反对者和\Quote{健全教义}\BibleRef{弟前 1:10}的仇敌的整体目标，就是通过夷平宗徒传统并彻底摧毁它，来动摇基督信德的基础。所以如同债务人——当然是\Emph{善意的}债务人——他们吵着要书面证据，并拒绝教父们未成文的传统，认为毫无价值。但我们不会松懈对真理的辩护。我们不会懦弱地放弃事业。主已经交付给我们一个必要而救恩的道理：圣神应与圣父并列。我们的反对者想法不同，他们觉得合适的分裂、撕裂，并将祂贬低为服役之神的本性。那么他们岂不是无可争辩地使自己的亵渎比主所规定的法律更具权威吗？来吧，放下单纯的争论。让我们按如下方式考虑面前的问题：\par

26. 我们成为基督徒是从何而来的？通过我们的信德，大家都会这样回答。而我们又是如何得救的呢？显然因为我们借着在圣洗中赐予的恩宠得以重生。还能有其他方式吗？在认识到这救恩是通过圣父、圣子和圣神而确立之后，我们竟要抛弃我们所接受的\Quote{那种教义形式}\BibleRef{罗 6:17}吗？如果我们现在发现离救恩比\Quote{初信之时}更远，并否认我们当时所接受的，岂不是该大大哀叹？无论一个人是未领圣洗就离世，还是领受了缺少传统某些要求的圣洗，他的损失是一样的。任何人若不能始终并处处坚持并紧握我们初入门时所记录的信仰宣认作为一种确切的保护（当时我们\Quote{从偶像中}，来到了\Quote{永生天主}面前\BibleRef{得前 1:9}），他就使自己成为天主\Quote{应许}\BibleRef{弗 2:12}的\Quote{局外人}，与自己亲笔所写的信德宣誓为敌。因为如果对我来说，我的圣洗是生命的开始，那重生的日子是第一天，那么显然在义子的恩宠中所说出的宣认是最尊贵的。难道我竟能被这些人的花言巧语所迷惑，而抛弃那引导我走向光明、赐给我认识天主的恩赐、使我这个因罪而成为仇敌的人成为天主的儿女的传统吗？至于我，我祈祷带着这个宣认离世去见主，我吩咐他们持守信德直到基督的日子，并保持圣神与圣父和圣子不可分割，在信仰宣认和荣耀颂中都保守他们在圣洗时所受的教导。\par

% structure: p0071 source=h2 target=section reason=relative-heading-rank; base=h2

\section{第十一章}

\Emph{否认圣神的人是背命者。}\par

27. \Quote{谁有祸患？谁有忧愁？}\BibleRef{箴 23:29} 谁有困苦和黑暗？谁有永罚？岂不是背命者吗？岂不是否认信德的人吗？他们否认的凭据是什么？岂不是他们轻视了自己的宣认吗？他们在何时宣认了什么？信仰圣父、圣子及圣神，当他们弃绝魔鬼及其天使，并说出那些救恩的话时。那么，光明之子该用什么合适的称呼来称呼他们？他们难道不是被称为背命者，因为他们违背了自己救恩的盟约？我该如何称呼否认天主的行为？什么是否认基督？岂不是背命？对于否认圣神的人，你希望我用什么称号？岂不是同样，因为他破坏了与天主的盟约？既然对圣神的信仰宣认确保真宗教的福分，而否认它将人置于不敬神的厄运之下，那么他们轻视宣认，不是出于对火、刀剑、十字架、鞭笞、轮刑或刑架的恐惧，而仅仅是被\OriginalTerm{圣神反对者}{pneumatomachi}的诡辩和诱惑所迷惑，这岂不可怕？我向每一个宣认基督却否认天主的人作证：基督对他毫无益处；向每一个呼求天主却拒绝圣子的人：他的信德是虚妄的；向每一个弃绝圣神的人：他对圣父和圣子的信德是无用的，因为没有圣神的临在，他甚至不能持守信德。因为不信圣神的人，不信圣子；不信圣子的人，不信圣父。因为除非借着圣神，没有人能\Quote{说耶稣是主}\BibleRef{格前 12:3}，并且\Quote{从来没有人看见过天主，只有那在父怀里的独生子神，祂将其显现出来。}\par

这样的人在真崇拜中无分无份；因为除非借着圣神，不能崇拜圣子；除非借着义子的圣神，不能呼求圣父。\par

% structure: p0075 source=h2 target=section reason=relative-heading-rank; base=h2

\section{第十二章}

\Emph{反对那些声称单以圣父之名的圣洗就够了的人。}\par

28. 愿无人因宗徒在提及圣洗时常常省略圣父和圣神之名而受误导，或因此以为称呼圣名不被遵守。他说：\Quote{你们众人，凡受洗归入基督的，都穿上了基督}；又说：\Quote{你们众人凡受洗归入基督的，正是受洗归入祂的死}。因为称呼基督就是对整体的宣认，由此显明了赐予的天主、接受圣子以及作为傅油的圣神。我们从宗徒大事录中从伯多禄学到：\BibleQuote{纳匝肋人耶稣，即天主以圣神傅了祂；}\BibleRef{宗 10:38}在依撒意亚中：\BibleQuote{主的神在我身上，因为主用油傅了我；}\BibleRef{依 60:1}圣咏作者说：\BibleQuote{因此天主，你的天主，以喜乐之油傅了你，胜过你的同伴。}然而，圣经在论及圣洗时，有时单独明确提及圣神。\par

因为经上说：\Quote{我们众人都在一个圣神内受洗，成了一个身体。}与此一致的是这些经文：\BibleQuote{你们要受圣神的洗}\BibleRef{宗 1:5}以及\BibleQuote{祂要用圣神给你们洗}\BibleRef{路 3:16}。但无人可因此称那仅呼求圣神之名的洗礼为圆满的洗礼。因为那藉赋予生命的恩宠所传授给我们的传统必须永远保持完整。那救赎我们生命于败坏的主给了我们更新的权能，其原因无法言说且隐藏于奥迹中，却给我们的灵魂带来伟大的救恩，以至于增减任何事物都明显意味着永生的失落。如果在圣洗中，将圣神与圣父及圣子分离对施洗者有害，对受洗者无益，那么我们与圣神、圣父和圣子的分离又如何能安全呢？信德和圣洗是两种相近且不可分的救恩之路：信德藉圣洗而圆满，圣洗因信德而坚固，两者皆藉同一圣名完成。因为正如我们信于圣父、圣子及圣神，我们也受洗于圣父、圣子及圣神之名；先是引导我们进入救恩的宣认，随后是圣洗，为我们赞同盖印。\par

% structure: p0079 source=h2 target=section reason=relative-heading-rank; base=h2

\section{第十三章}

\Emph{陈述保禄书信中天使何以与圣父和子并列的理由。}\par

29. 然而，有人反对说，其他受造物虽与圣父和子一同被列举，却并非总是同受荣耀。例如，宗徒在给弟茂德的训命中，以这样的话将天使与他们并列：\BibleQuote{我在天主和基督耶稣及蒙选的天使面前嘱咐你}\BibleRef{弟前 5:21}。我们并不将天使从其余受造物中分离，然而有人争辩说，我们也不容许他们与圣父和子同列。对此我答复说，尽管这个论点因其明显荒谬而不值得回答，但在一位温和公正的审判官面前，尤其是一位以宽待受审者来彰显其判决无可指摘的公义者面前，一个人或许愿意让同是奴隶的同伴作证；但使奴隶得自由、被称为天主的子女并脱离死亡而活，只能由那与我们具有自然亲属关系、并从奴隶地位转变的那一位来成就。因为一个外人如何能使我们与天主成为亲属？一个自身仍处奴役之下者如何能释放我们？由此可见，提及圣神与提及天使并非在同等条件下。圣神被称为生命之主，而天使则是同为奴隶者的同盟和真理的忠实见证。圣人们惯常在证人面前传达天主的诫命，如宗徒亲自对弟茂德说：\BibleQuote{你在许多证人面前从我听见的，你要把这些交托给忠信的人}\BibleRef{弟后 2:2}；如今他召天使作证，因为他知道当主在祂父的荣耀中来临，以公义审判世界时，天使将同主在一起。因为祂说：\BibleQuote{凡在人面前承认我的，人子也要在天主的众天使面前承认他；凡在人面前否认我的，必在天主的众天使面前被否认}\BibleRef{路 12:8}；保禄在另一处说：\BibleQuote{那时主耶稣从天降来，同祂的众天使}\BibleRef{得后 1:7}。因此他已在天使面前作证，为伟大的审判台预备了美好的证据。\par

30. 不仅保禄，凡是受托任何言语职务的人，都不断地作证，呼天唤地作见证，因为现今所做的一切事都在它们内进行，在审查生命的一切行为时，它们将与被审判者一同在场。所以经上说：\Quote{他向高天和大地呼唤，为审判他的子民。}梅瑟即将向百姓传达神谕时说：\BibleQuote{我今日呼天唤地向你们作证；}\BibleRef{申 4:26}又在歌中说：\BibleQuote{诸天哪，倾听吧，我要发言；大地啊，聆听我口中的言语；}\BibleRef{申 32:1}依撒意亚：\BibleQuote{诸天哪，倾听吧；大地啊，侧耳听；}\BibleRef{依 1:2}耶肋米亚描写天听闻百姓不圣洁的行为而惊骇：\Quote{天为此惊异，极其恐惧，因为我的人民犯了两样恶事。}同样，宗徒知道天使被指派为人师表和保护者，便呼召他们作证。此外，农的儿子若苏厄甚至立了一块石头作他话语的见证（雅各伯曾称一堆石堆为见证），\BibleRef{创 31:47}因为他说：\BibleQuote{看，这石头今天要作你们的见证，直到末日的日子，当你们欺骗上主我们的天主时。}也许是相信藉天主的能力，连石头也会说话来定犯法者的罪；或者即使不然，至少每个人的良心会因提醒的力量而受创伤。这样，受托管理灵魂的人提供任何可能的见证，以便在将来某日出示它们。但是圣神与天主同列，不是因为一时的急需，而是因为本性的共融；不是被我们强拉进来，而是被主所邀请。\par

% structure: p0083 source=h2 target=section reason=relative-heading-rank; base=h2

\section{第十四章}

\Emph{有人曾受圣洗归于梅瑟并相信他，对此的反对意见及答复；兼论预像。}\par

31. 但有人辩称：即便有人受圣洗归于圣神，也不能因此将圣神与天主同列。有人\BibleQuote{曾在云中和海里受圣洗归于梅瑟。}\BibleRef{格前 10:2}而且承认信德以往也曾寄托于人，因为\Quote{百姓信了天主和他的仆人梅瑟。}那么，有人问道：既然有证据表明同样的事以前也曾对人说过，我们为何因信德和圣洗将圣神高举远在受造物之上？那么，我们该如何回答？我们的回答是：对圣神的信德与对圣父和圣子的信德相同；同样，圣洗也是如此。但对梅瑟和云彩的信德，如同在影子和预像中。神性常藉预像的粗糙模糊轮廓来表现；但因为神圣事物由微小的人间事物所预表，显然我们不能因此断定神性是微小的。预像是所期待之事的展示，并给予未来的模仿性预尝。因此亚当是\BibleQuote{那要来的那一位}\BibleRef{罗 5:14}的预像。在预像中，\BibleQuote{那盘石就是基督；}\BibleRef{格前 10:4}水是圣言活力的预像；正如他说：\BibleQuote{谁若渴，让他到我这里来喝。}\BibleRef{若 7:37}玛纳是从天降下的生命之粮的预像；\BibleRef{若 6:49}旗杆上的铜蛇是藉十字架完成的救恩苦难的预像，因此连仰望它的人都得保存。同样，以色列出埃及的历史被记载，以显示那些藉圣洗而得救的人。因为以色列的长子被保存，如同受圣洗者的身体，因给那些用血标记的人施予恩宠而得保存。因为羔羊的血是基督血之预像；长子是首先成形者的预像。既然首先成形者必然存在在我们内，并藉代代相传直至终结，那么我们在\BibleQuote{亚当内众人都死了}\BibleRef{格前 15:22}，并且\BibleQuote{死亡为王}\BibleRef{罗 5:17}直到法律成全和基督来临。长子被天主保存免受毁灭者触摸，以显示我们在基督内成为活人，不再在亚当内死去。海洋和云彩在暂时带领人藉惊奇走向信德，但为将来，它们预表将要来临的恩宠。\BibleQuote{谁明智，他会明白这些事？}\BibleRef{欧 14:9}——海洋如何预表一种带来法郎脱离的圣洗，同样，这洗洁带来魔鬼暴政的脱离。海洋在其内杀死了敌人：在圣洗中，我们对天主的敌意也死去。百姓从海洋安然走出：我们也好似从死者中复活，从水中上来，\BibleQuote{得救，是藉那召叫我们的那一位的恩宠。}\BibleRef{弗 2:5}云彩是圣神恩赐的阴影，圣神藉我们\BibleQuote{死我们的肢体}\BibleRef{哥 3:5}冷却我们情欲的火焰。\par

32. 那么，怎样呢？因为他们曾象征性地受洗归于梅瑟，难道圣洗的恩宠就因此微小吗？若是如此，若我们在每一情形下都以类型来贬损我们特恩的尊荣，那么连这些特恩中的任何一个都不能算作伟大；那么天主为我们的罪赐下祂独生子的爱，就不是伟大而超凡的，因为亚巴郎没有吝惜自己的儿子；那么主的苦难也就不是光荣的，因为一只羊预表了代替依撒格的祭献；那么下降地狱并不令人畏惧，因为约纳先前已预表了三日三夜的死亡。所有以影子判断实体，并将预表者与类型相比较，试图通过梅瑟和海来同时贬抑整个福音救世计划的人，都在圣洗之事上作了同样有损的比较。在海中有什么罪的赦免、生命的更新呢？通过梅瑟有什么属灵的恩赐呢？有什么罪的死亡呢？那些人没有与基督同死，所以他们也没有与祂同活。他们没有\BibleQuote{带着属天的肖像}\BibleRef{格前 15:49}，他们没有\BibleQuote{在身上带着耶稣的死状}\BibleRef{格后 4:10}，他们没有\BibleQuote{脱去旧人}，他们没有\BibleQuote{穿上了新人，这新人是在知识上按着创造他者的肖像而更新的}\BibleRef{哥 3:9}。那么，你为什么将那些只在名称上相同，而事物本身之间的差别却如同梦境与现实、影子与形像的实体之间的差别一样大的洗礼相提并论呢？\par

33. 但是，对梅瑟的信仰不仅不表明我们对圣神的信仰毫无价值，而且，如果我们采用对手的论证路线，它反而削弱我们对宇宙天主的宣认。经上记载：\BibleQuote{百姓信服了天主和祂的仆人梅瑟}\BibleRef{出 14:31}。梅瑟于是与天主联结，而非与圣神；他预表的不是圣神，而是基督。因为那时，在法律的服务中，他藉着自己预表了\BibleQuote{天主与人之间的中保}\BibleRef{弟前 2:5}。梅瑟在关乎天主的事上为人民调解时，并非圣神的仆人；因为法律是\BibleQuote{藉着天使，经中保之手而设立的}\BibleRef{迦 3:19}，即梅瑟，应百姓的请求：\BibleQuote{你同我们说话……不要让天主同我们说话}\BibleRef{出 20:19}。因此，对梅瑟的信仰归于主，即天主与人之间的中保，祂曾说：\BibleQuote{你们若信梅瑟，也必信我}\BibleRef{若 5:46}。那么，我们对主的信仰是小事吗，因为它曾预先通过梅瑟被指明？所以，即使有人曾受洗归于梅瑟，也不意味着在圣洗中由圣神所赐的恩宠是微小的。我也可以指出，圣经中通常将梅瑟和法律并提，如经上说：\BibleQuote{他们有梅瑟及先知}\BibleRef{路 16:29}。因此，当要提及法律的洗礼时，经文说：\BibleQuote{他们受洗归于梅瑟}\BibleRef{格前 10:2}。那么，这些真理的诽谤者为何要通过影子和类型，企图藐视和嘲讽我们\BibleQuote{盼望的喜悦}\BibleRef{希 3:6}以及我们的天主和救主的丰厚恩赐呢？祂藉着重生使我们的青春如鹰一般更新。这实在是极其幼稚，如同一个仍需吃奶的婴孩，不明白我们救恩的伟大奥迹；因为按照我们教育的渐进过程，当我们在虔敬的训练中被引向完美时，我们首先被教导与我们的智力相适应的初阶和较容易的功课，而分配我们命运的那一位，总是逐渐引导我们，如同将习惯于黑暗的眼睛逐渐适应真理的大光。祂顾惜我们的软弱，并在祂智慧丰富的深处\BibleRef{罗 11:33}和祂聪慧的不可测度的判断中，使用了这种适合我们需要、温和的对待，逐渐使我们习惯于先看事物的影子，并观看水中的太阳，以免我们撞上纯粹无杂的光而失明。同样，法律拥有未来之事的影子，而先知们的典型教导是对真理的晦涩表达，这些都是操练心灵之眼的工具，以便由此过渡到隐藏在奥迹中的智慧\BibleRef{格前 2:7}变得容易。关于预表，至此已足够；事实上，也不可能在此话题上逗留更久，否则附带讨论会比主要论证庞大许多倍。\par

% structure: p0088 source=h2 target=section reason=relative-heading-rank; base=h2

\section{第十五章}

答复所提的异议：我们受洗是\Quote{归于水}。亦论及圣洗圣事。\par

34. 还有什么呢？的确，我们的对手准备充分。他们争辩说：我们受洗归于水，当然我们不会尊崇水超过一切受造物，或将水的尊荣与圣父和圣子分享。这些人的论证如同愤怒的辩论者所常做的，在攻击冒犯他们的人时无所不用其极，因为他们的理性被情绪所蒙蔽。然而，我们不会回避对这些问题的讨论。即使我们教导不了无知者，至少我们在恶者面前不会退缩。但让我们稍作回溯。\par

35. 我们的天主和救主关于人的经济，是使人从堕落中被召回，并从因不顺从而导致的与天主相疏远中返回，以达至与天主密切的共融。这就是基督降生成人、福音中所描述的模范生活、苦难、十字架、坟墓和复活的原因；这样，通过效法基督而得救的人，便重新获得了那旧时的义子地位。为达至生命的成全，效法基督是必要的，不仅效法他在生活中为我们树立的温良、谦卑和忍耐的榜样，也要效法他实际的死亡。因此，效法基督的保禄说：\BibleQuote{使我相似他的死亡，以便我达到死人的复活。}\BibleRef{斐 3:10} 那么，我们如何相似他的死亡呢？\BibleRef{罗 6:4} 因着圣洗，我们与他同葬了。那么，这种埋葬的方式是什么？效法所带来的益处又是什么？首先，必须切断旧生命的连续性。除非人重生，否则这是不可能的，正如主的话：\BibleRef{若 3:3} 因为重生，正如这名称所示，是第二生命的开始。因此，在开始第二生命之前，必须结束第一生命。正如赛跑者转身开始第二段赛程时，在相反方向的运动之间会有一种停顿和间歇；同样，在改变生命时，死亡作为两者之间的中介似乎是必要的，它结束一切先前的事，开始一切后来的事。那么，我们如何达成下降阴府呢？通过圣洗效法基督的埋葬。因为受洗者的身体仿佛被埋葬在水中。因此，圣洗象征性地表示除去肉体的作为；正如宗徒所说：你们\BibleQuote{受了不是人手所行的割损，而是藉基督的割损，脱去肉欲之身的割损；与他同葬于圣洗中。}\BibleRef{哥 2:11} 此外，灵魂得以洗涤，除去因肉性心思而沾染的污秽，正如经上记载：\Quote{你洗涤我，我就比雪更白。} 因此，我们不像犹太人那样每次沾染污秽就洗涤，而是承认只有一个救恩的圣洗。因为为世界的死亡只有一次，死人的复活也只有一次，圣洗正是这预象。为此，主——我们生命的施予者——赐给我们圣洗的盟约，其中包含生命与死亡的预象，因为水实现了死亡的图像，而圣神赐给我们生命的保证。由此，我们问题的答案便清楚了：为何水与圣神相关联？原因是圣洗有两个目的：一方面，消灭罪身，使它不再结出死亡的果实；另一方面，我们为圣神而生活，使果实成圣；水如同坟墓般接纳身体，象征死亡，而圣神注入赋予生命的能力，使我们的灵魂从罪的死亡中更新，恢复原初的生命。这就是从水和圣神重生的意义：死亡在水里完成，而生命藉圣神在我们内运作。这样，三次浸入和三次呼求，圣洗的伟大奥迹得以完成，为的是完全象征死亡的预象，并使受洗者因神圣知识的传授而灵魂得到光照。因此，若水中有什么恩宠，那不是出于水的本性，而是出于圣神的临在。因为圣洗\BibleQuote{并不是除去肉体的污秽，而是向天主求一纯洁的良心。}\BibleRef{伯前 3:21} 所以，为训练我们迎接复活后的生命，主规定了福音所要求的一切生活方式，为我们制定了温良、忍受不公、摆脱贪图享乐所带来的污秽和贪欲的律法，为使我们能刻意地预先赢得并实现来生本然具有的一切。因此，若有人试图给福音下定义，说它是复活后生命的预兆，我认为这并不失当。现在让我们回到主要议题。\par

36. 藉着圣神，我们恢复乐园，升入天国，重获义子名分，得以自由地呼求天主为父，成为\Person{基督}恩宠的分享者，被称为光明之子，分享永生的荣耀，总而言之，我们被带入一切\Quote{圆满的祝福,}的状态\BibleRef{罗 15:29}，无论在现世还是来世，对于为我们存留的一切美好恩赐，藉着应许，藉着信德，我们看见它们恩宠的返照，如同已临在，而等待着圆满的享有。如果这是预尝，那么完美是什么？如果这是初果，那么完全的实现是什么？此外，由此也可以领会来自圣神的恩宠与水的洗礼之间的区别：因为\Person{若翰}确实以水施洗，但我们的主\Person{耶稣基督}以圣神施洗。他说：\BibleQuote{我固然用水洗你们，为使你们悔改；但在我以后来的那一位，比我更强，我连替祂提鞋也不配；祂要以圣神和火洗你们。}\BibleRef{玛 3:11}这里祂称审判时的考验为火的洗礼，正如宗徒所说：\BibleQuote{火要试验各人的工程，看它是哪一种。}\BibleRef{格前 3:13}又说：\BibleQuote{那日子要表明它，因为它是藉着火显现的。}\BibleRef{格前 3:13}而且，以前也有一些人为维护真宗教，为了\Person{基督}而殉道，不是仅仅象征性地，而是实际地，因此他们不需要任何外在的水的记号来获得救恩，因为他们是用自己的血受洗的。我这样说并非贬低水的洗礼，而是要推翻那些自高反对圣神的论点；他们混淆了彼此不同的事物，比较那些无法比较的事物。\par

% structure: p0093 source=h2 target=section reason=relative-heading-rank; base=h2

\section{第十六章}

\Emph{圣神在任何概念中都与圣父和圣子可分离，无论是在可感知之物的创造中，在人类事务的安排中，还是在将来的审判中。}\par

37. 让我们回到最初提出的论点：在一切事物中，圣神与圣父和圣子都是不可分离的，完全不能被分开。\Person{圣保禄}在关于说方言之恩赐的段落中，写信给格林多人说：\BibleQuote{如果你们都说预言，有一个不信的人或一个未受教的人进来，他被众人所责备，被众人所审断，他心中的隐密就显露出来；于是他会俯伏在地，朝拜天主，并宣布天主真是在你们当中。}\BibleRef{格前 14:24}那么，如果天主藉着说预言（即按照圣神恩赐的分施而行动）而被知道在先知们当中，就让我们的反对者思考，他们将把什么地位归于圣神。让他们说，是把祂位列与天主同等的地位更合适，还是把祂推到受造物的位置？\Person{伯多禄}对\Person{撒斐辣}说的话：\BibleQuote{你们怎么同心试探主的圣神？你不是欺骗人，而是欺骗天主。}表明得罪圣神与得罪天主是一样的；这样你可以学到，在每一个行动中，圣神都与圣父和圣子紧密相连，不可分离。天主实行不同的功效，主实行不同的职分，但与此同时，圣神也按照自己的旨意临在，按照每个领受者的功过分施恩赐。因为经上说：\BibleQuote{神恩虽有区别，却是同一的圣神；职分虽有区别，却是同一的主；功效虽有区别，却是同一的天主，在众人身上行一切。}\BibleQuote{这一切都是这同一的圣神所行的，随祂的意愿，个别地分赐给每个人。}\BibleRef{格前 12:11}然而，不可因为在此段经文中宗徒先提到圣神，其次提到圣子，再其次提到天主圣父，就认为他们的等级被颠倒了。宗徒只是按照我们的思维习惯开始；因为当我们领受恩赐时，首先出现在我们面前的是分施者，然后我们想到派遣者，接着我们将思想提升到恩惠的泉源和原因。\par

38. 此外，从起初所造的万物可以学到圣神与圣父及圣子的共融。纯洁、理智且超世的掌权者之所以被称为圣，是因为他们从由圣神赐予的恩宠中获得圣德。因此，天上诸能的受造方式被略过不提，因为宇宙起源的叙述者只向我们揭示了感官可感知之物的创造。但是，你既能从可见之物形成类比以理解不可见之物，就当光荣那创造万有的造物主，凡可见与不可见之物、率领者与掌权者、权威、宝座与宰制者，以及所有其他我们无法称呼的理性本性，都是祂所造的。在创造中，我请你首先思考万受造之物的原始因——圣父；创造因——圣子；成全因——圣神；如此，服役的神灵因圣父的旨意而存立，因圣子的运作而受造，因圣神的临在而得以成全。此外，天使的完美在于圣化并持守于圣化中。但愿没有人想象我是在肯定有三个原初的\OriginalTerm{位格}{hypostases}，或是在暗示圣子的运作是不完美的。因为存在之物的第一原理是独一的，祂通过圣子创造，并通过圣神成全。那在万有中运作万有的圣父，其运作并非不完美；同样，圣子的创造之工若未经圣神的成全，也不是不完全的。圣父单凭自己的旨意创造，不需要圣子，然而祂愿意通过圣子创造；同样，圣子依照圣父的样式运作，本来无需合作，但圣子也愿意通过圣神来成全。因为\BibleQuote{诸天因上主的一句话而创造，万军因祂口中的嘘气（圣神）而形成。}那么，圣言并非仅仅是通过发音器官承载于空气上的有意义印记；祂口中的嘘气也不是通过呼吸器官呼出的雾气；但圣言就是那\BibleQuote{在起初就与天主同在}并且\BibleQuote{就是天主}的一位，\BibleRef{若 1:1}；而天主口中的嘘气就是\BibleQuote{那由父发出的真理之神}，\BibleRef{若 15:26}。因此，你当体认三者：那位发出命令的主，创造万有的圣言，以及使之坚定的圣神。而坚定，若非按照圣德的成全，又能是什么呢？这成全表达了坚定的牢固、不变和固着于善。但没有圣神就没有圣化。天上诸能并非本性为圣；若是如此，在这方面他们与圣神之间就没有区别。他们之获得圣德，是来自圣神，按他们各自相对于圣神的卓越程度。烙铁与火一同被认识；然而材质与火是截然不同的。同样，天上诸能也是如此：他们的本体或许是气态的灵，或无质料的火，正如经上所载：\BibleQuote{祂使祂的天使成为神灵，祂的臣仆成为火焰}；因此他们存在于空间中，并成为可见的，以适当的形体显现给配得的人。但他们的圣化，由于外在于他们的本体，通过圣神的相通而促成他们的成全。他们因持守于善与真而保持其等级，并且当他们保留意志的自由时，他们从不从对那位真正之善的坚忍侍奉中坠落。其结果是，若你以你的论证废除了圣神，天使的万军就会解散，总领天使的宰制就会被毁灭，一切陷入混乱，他们的生命失去法律、秩序和分别。因为天使们若不藉圣神的德能，怎能呼喊\BibleQuote{在天受光荣于天主}？\BibleRef{路 2:14} 因为\BibleQuote{没有人能说\InnerQuote{耶稣是主}，除非藉圣神；也没有人受天主圣神感动而说\InnerQuote{耶稣是可诅咒的}}；\BibleRef{格前 12:3} 这可能为邪恶敌对的神灵所说，他们的堕落证实了我们关于不可见掌权者意志自由的论述；因为他们处于善恶之间的平衡状态，并因此需要圣神的援助。我确实认为，连\Person{加俾额尔}\BibleRef{路 1:11}也并非以其他方式预言将来的事，而是藉圣神的预知，因为圣神所分施的恩赐之一是预言。那位被派去向蒙爱的人宣告神视奥秘的人，他所获得的启示隐藏之事的智慧，若不是来自圣神，又能从哪里得来呢？启示奥秘确实是圣神的独特功能，正如经上所载：\BibleQuote{天主藉圣神将这一切启示给我们}。\BibleRef{格前 2:10} 并且，\BibleQuote{宝座、宰制者、率领者与掌权者}\BibleRef{哥 1:16}若不\BibleQuote{瞻仰那在天之父的面容}\BibleRef{玛 18:10}，又怎能过他们幸福的生活呢？但没有圣神就不可能瞻仰！正如在夜间，若你从屋内撤去灯光，眼睛就会失明，其官能也失去活力，无法辨别物体的价值，金子因无知而如铁一般被践踏；同样，在理智世界的秩序中，没有圣神，律法的高尚生活就无法存续。因为如此存续的可能性，如同军队在没有指挥官时维持纪律，或合唱团在没有领唱者引导下保持和谐一样。色辣芬怎能呼喊\BibleQuote{圣、圣、圣}？\BibleRef{依 6:3} 除非他们受圣神教导，真实宗教要求他们多少次在赞美中提高声音？难道\Quote{祂的众天使}和\Quote{祂的万军}赞美天主吗？这是藉圣神的合作。难道\BibleQuote{一千千}的天使侍立在祂面前，以及\BibleQuote{一万万}的服役神灵？\BibleRef{达 7:10} 他们藉圣神的力量无可指责地做他们适当的工作。至高天上一切光荣而不可言喻的和谐，无论是在事奉天主方面，还是在天上诸能的相互和谐中，都只能藉圣神的引导得以保存。因此，对于那些不是通过增长和进步逐渐成全，而是从受造之初就是完美的存在，在受造中就有圣神的临在，祂从自己那里赐予他们涌流的恩宠，以完成和完美他们的本质。\par

39. 然而，当我们论及我们的伟大天主和救主耶稣基督为人类所行的救世计划时，谁能否认这些计划皆藉圣神的恩宠完成呢？无论你愿意查考古时的证据——圣祖们的祝福、藉法律所赐的援助、预象、预言、战争中的英勇事迹、藉义人所行的奇事；——或者另一方面，在主降生成人的救世计划中所行的一切——一切都是藉着圣神。首先，祂被傅油，并且不可分离地临在于主的肉身中，正如经上所载：\BibleQuote{你看见圣神降下，停在谁身上，谁就是}\BibleRef{若 1:33}\BibleQuote{我的爱子}\BibleRef{玛 3:17}，以及\BibleQuote{纳匝肋人耶稣}，\BibleQuote{天主以圣神傅了祂}\BibleRef{宗 10:38}。此后，每一作为都是在圣神的合作下完成的。当主受魔鬼试探时，圣神在场；因为经上说：\BibleQuote{耶稣被圣神领往旷野，受魔鬼试探}\BibleRef{玛 4:1}。祂在行奇事时与主不可分离；因为经上说：\BibleQuote{我若仗赖天主的神驱魔}\BibleRef{玛 12:28}。祂在主从死者中复活后也没有离开祂；因为当主更新人，并向门徒们嘘气，恢复人因天主嘘气而失落的恩宠时，主说了什么？\BibleQuote{你们领受圣神吧！你们赦免谁的罪，就给谁赦免；你们留存谁的，就给谁留存}\BibleRef{若 20:22}。教会的秩序是藉圣神建立的，这岂不是明显而无可争辩的吗？因为经上说：\BibleQuote{在教会内，第一是宗徒，第二是先知，第三是教师，其次是异能，再次是治病的恩赐、助人的、治理的、各种语言的}\BibleRef{格前 12:28}，因为这秩序是按圣神恩赐的分赐而设立的。\par

40. 此外，凡善用理智的人都会发现，即使在主自天而来的期待显现的时刻，圣神也并非如某些人所猜想的那样无事可做；相反，甚至在祂显现的日子，那位\BibleQuote{那真福的、惟一全能者}\BibleRef{弟前 6:15}\BibleQuote{将以正义审判天下}\BibleRef{宗 17:31}时，圣神也要与祂同在。因为谁会对天主为义人所预备的美善如此无知，竟不知道义人的冠冕就是圣神的恩宠，在那一天将以更丰盛、更完美的尺度赐予，那时属灵的荣耀将按各人英勇作战的程度分赐呢？因为在圣人们的荣耀中，父家中有\Quote{许多住处}，即尊荣的差异：因为\BibleQuote{这星和那星的荣耀原有分别；死人的复活也是这样}\BibleRef{格前 15:41}。因此，那些受了圣神的印记，以待得救赎的日子的人，并保持所领受的圣神初果纯洁无损的人，将听见这些话：\BibleQuote{好！善良忠信的仆人，你既在少许事上忠信，我必委派你管理许多大事}\BibleRef{玛 25:21}。同样，那些因自己行为的邪恶而使圣神忧郁的人，或没有为那赐予者而工作的人，将丧失他们所领受的恩宠，其恩宠将转给他人；或按一位福音作者的话，他们甚至要被\BibleQuote{斩成两半}\BibleRef{玛 24:51}——这\Quote{斩成两半}意指与圣神完全分离。身体不会被分开，一部分受罚，另一部分幸免；因为若整体犯了罪，只有一半受罚，那就像古老的寓言，有失公正的审判者。灵魂也不会被切成两半——那整个灵魂都充满罪恶情感，并协同身体行恶。正如我所指出的，\Quote{斩成两半}就是灵魂与圣神永远分离。因为如今，尽管圣神不与不配的人相混，但祂似乎仍以某种方式临在于那些曾受印记的人，等待他们悔改后的得救；但那时，祂将完全从玷污祂恩宠的灵魂上被割断。因此\Quote{阴府中没有称谢者；死亡中没有纪念天主的}，因为圣神的帮助不再存在。那么，怎能设想审判能在没有圣神的情况下进行呢？因为经上指出，祂本身就是义人的赏报\BibleRef{斐 3:14}，那时以完美代替抵押，而罪人的初次定罪，就是被剥夺他们似乎拥有的一切。但圣神与圣父及圣子联合的最大证明，就是祂与天主的关系，如同我们内的灵与我们每个人的关系。因为经上说：\BibleQuote{除了人内里的灵，谁知道人的事？同样，除了天主的圣神，也没有人知道天主的事}\BibleRef{格前 2:11}。\par

关于这一点，我已说得足够了。\par

% structure: p0100 source=h2 target=section reason=relative-heading-rank; base=h2

\section{第十七章}

\Emph{反对那些说圣神不应与圣父及圣子同列，而应列于其下的人。其中也对有关正当位列的信仰作了简要说明。}\par

41. 然而，他们所谓的\OriginalTerm{次第计数}{sub-numeration}，以及他们使用这个词的含义，若不费一番心思，甚至无法想象。众所周知，这个词是从\Quote{世俗的智慧}引入我们语言的；但我们现在要考虑的是，它是否与当前讨论的主题有直接关系。那些精通无用研究的人告诉我们，有些名词是普遍的，外延广泛，另一些则更具体，有些的力量比另一些更受限制。例如，本质是一个普遍名词，可断言于一切有生命和无生命的事物；而动物则更具体，可断言的对象比前者少，但比其下位者多，因为它包含理性与非理性本性。再者，人比动物更具体，而个体的人比人更具体，比个体的人更具体的是\Person{伯多禄}、\Person{保禄}或\Person{若望}。那么，他们是否通过次第计数指将普遍者划分为其从属部分？但我难以相信他们会痴迷到如此程度，竟断言宇宙的天主，如同某种仅凭理性可构想、在任何\OriginalTerm{位格}{hypostasis}中并无实际存在的普遍性质，被划分为从属部分，然后这种划分被称为次第计数。即使精神忧郁的疯狂之人也不会说这样的话，因为除了不虔敬之外，他们正在建立与自己论点完全相反的论证。因为从属部分与它们从中划分出来的部分本质相同。这种荒谬的明显性使我们难以找到论据来反驳他们的不合理；因此，他们的愚昧实际上看起来像是一种优势；就像柔软而有弹性的物体无法承受重击一样。不可能对明显的荒谬进行有力的反驳。我们唯一的选择是默默地绕过他们可憎的不虔敬。然而，我们对弟兄的爱和对手的纠缠不休使我们无法保持沉默。\par

42. 他们主张什么？看看他们欺骗的言辞。\Quote{我们断言，\OriginalTerm{同列计数}{connumeration}适用于同等尊严的主体，而\OriginalTerm{次第计数}{sub-numeration}适用于那些向低劣方向变异的主体。} \Quote{为什么？} 我回答说，你们这样说？我无法理解你们非凡的智慧。你们是否认为黄金与黄金同列计数，而铅因材料的低贱不配同列计数，而是被次第计数于黄金？你们是否赋予数字如此大的重要性，以至于它能提升低贱之物的价值，或摧毁贵重之物的尊严？因此，你们又会让黄金次第计数于宝石之下，让那些较小且无光泽的宝石次第计数于较大且颜色更亮的宝石之下。但是，那些除了\BibleQuote{讲述或听一些新奇的事}\BibleRef{宗 17:21}以外不做任何事的人，还有什么不会说出来呢？让这些支持不虔敬的人今后与斯多葛派和伊壁鸠鲁派同列吧。对于价值较低的事物与价值极高的事物之间，怎么可能进行次第计数？一个铜币如何次第计数于一个金币？\Quote{因为，} 他们回答说，\Quote{我们不说拥有两个钱币，而是一个和一个。} 但其中哪一个被次第计数于另一个？每个都被同样提及。如果你们分开计数每个，你们通过以同样方式计数使它们价值相等；但如果你们合并它们，你们通过一个与另一个一起计数使它们的价值合一。但如果次第计数属于第二个被计数的，那么计数者可以靠先数铜币来支配。然而，让我们略过对他们无知的驳斥，将我们的论证转向主要话题。\par

43. 你是否主张圣子被列于圣父之下，圣神被列于圣子之下，还是你只将\OriginalTerm{下级编号}{sub-numeration}限制于圣神自身？另一方面，如果你将这下级编号也应用于圣子，你就复兴了同样的不敬神学说——本质的不同、等级的卑微、后来的存在，并且，通过这一术语，你将再次清楚地使所有反对\OriginalTerm{独生子}{Only-begotten}的亵渎之言重蹈覆辙。反驳这些亵渎之言的工作超出了我当前目的所能容许的范围；而且我更不必承担此任，因为这些不敬神之处已在别处尽我所能被驳斥了。另一方面，如果他们以为下级编号仅有益于圣神，那么必须教导他们：圣神与主相提并论的方式，恰与圣子与圣父相提并论的方式相同。主教导我们\Quote{父、子、及圣神之名}\BibleRef{玛 28:19}以同样的方式传授，并且，根据洗礼中所传授的言语的并列关系，圣神对圣子的关系与圣子对圣父的关系相同。如果圣神与圣子并列，圣子与圣父并列，那么显然圣神也与圣父并列。当这些名称被排列在同一并列系列中时，还有什么余地一方面说\OriginalTerm{同级编号}{connumeration}，另一方面说下级编号？不，毫无例外，有什么东西因被编号而失去了自己的本性？难道不是事物在被编号时保持其自然和原始的状态，而编号在我们中间被采用作为指示主体多样性的符号吗？我们计数一些物体，测量一些，称重一些；那些本性连续的东西我们通过测量来掌握；那些分开的东西我们使用编号（除了那些因精细而被测量的）；而重物通过天平的倾斜来区分。然而，这并不意味着，因为我们为了方便发明了帮助我们达到数量知识的符号，我们就改变了所指事物的本性。我们不说\Quote{在下方称量}那些被称量的事物，即使一个是金子另一个是锡；也不说\Quote{在下方测量}那些被测量过的事物；同样，我们也不会说\Quote{在下方编号}那些被编号的事物。如果其他事物都不容下级编号，他们怎么能声称圣神应当被下级编号呢？由于他们受异教谬误的困扰，他们想象那些在等级上或在本体的从属上处于低等的事物应当被下级编号。\par

% structure: p0105 source=h2 target=section reason=relative-heading-rank; base=h2

\section{第十八章}

\Emph{在承认三个\OriginalTerm{位格}{hypostases}时，我们如何保持\OriginalTerm{独一神论}{Monarchia}的敬神教义。其中也有对声称圣神被下级编号之人的反驳。}\par

44. 在传授父、子及圣神的名号时\BibleRef{玛 28:19}，我们的主并未将恩赐与编号相连。祂没有说\Quote{进入第一、第二和第三}，也没有说进入一、二、三，而是藉着神圣的名号，赐予我们认识那引人得救的信德的恩惠。因此，拯救我们的是我们的信德。编号被发明出来，作为指示对象数量的符号。但这些人，从每个可能的源头自取灭亡，甚至把计数的能力也转而反对信德。没有任何其他事物因编号的添加而经历任何变化，而这些人在神圣本性上却尊敬编号，以免他们超过对护慰者应有的尊敬之限度。但是，最智慧的先生们，让那不可接近者完全超越编号之上，正如希伯来人的古老敬畏用特殊的字符书写天主不可言说的名字，从而努力表达其无限的卓越。如果需要数算，你们可以数算；但不可因数算而损害信德。要么以沉默尊敬那不可言者；要么以符合真宗教的方式数算神圣事物。天主圣父是一位，\OriginalTerm{独生子}{Only-begotten}是一位，圣神是一位。我们个别地宣认每一个\OriginalTerm{位格}{hypostasis}；当我们必须数计时，我们不让无知的算术把我们带到众多天主的观念中。\par

45. 因为我们并非以叠加的方式计数，逐渐从单一增至众多，并说一、二、三——也不说第一、第二、第三。因为\Quote{我，}天主\Quote{是首先的，也是末后的。}\newline{}\BibleRef{依 44:6} 至今，甚至现在，我们从未听说过第二位天主。我们钦崇出自天主的天主，既承认位格的区别，同时也持守\OriginalTerm{独一统御}{Monarchy}。我们不将神学瓦解成分裂的众多，因为可以说，一种联合在神性不变中的\OriginalTerm{形态}{Form}，在天主圣父和天主独生子中被看见。因为子在天父内，天父也在子内；因为后者如何，前者也如何，而前者如何，后者也如何；这就是合一。因此，按照位格的区别，两者既是一个又是一个，而按照本性的共通，则是一个。那么，如果一个又是一个，为什么不是两位天主？因为我们谈论一位君王和君王的肖像，而不是两位君王。威严不被劈成两半，荣耀也不被分割。对统治和权柄的掌管是唯一的，因此我们归给的光荣颂也不是多个，而是一个；因为对肖像所表示的敬意传至原型。在一种情况下，肖像因模仿而如此，在另一种情况下，子因本性而如此；就如同在艺术作品中，相似性取决于形式，在神圣且单纯的本性中，合一在于神性的相通。此外，\OriginalTerm{圣神}{Holy Spirit}是唯一的，我们单独谈论祂，祂通过唯一子与唯一父相连，并通过祂自己完成可钦崇的、应受赞美的\OriginalTerm{天主圣三}{Trinity}。祂与父和子的密切关系，由祂不被列于受造物的众多之中，而是被单独谈论这一事实充分表明；因为祂不是众多中的一位，而是唯一者。因为如有一位父和一位子，同样也有一位圣神。因此，祂远离受造的本性，如同理性要求单一者远离复合的、众多的物体一样；祂与父和子合一，如同单一者与单一者相连。\par

46. 并且，我们关于本性共通的证明不仅来自这个源头，也来自祂被称为\Quote{出于天主}\newline{}\BibleRef{格后 1:12}的事实；当然不是像\Quote{万物都出于天主}那样的意义，而是指从天主而出，不是如同子那样通过生，而是如同祂口中的气息。但\Quote{口}绝不是肢体，圣神也不是会消散的气息；相反，\Quote{口}一词的使用仅限于它能适用于天主的程度，而圣神是一个具有生命、拥有至高圣化能力的实体。这样，密切的关系就显明了，而不可言喻存在的方式得到了保障。祂还被称作\Quote{基督的圣神}，因为祂在本质上与基督密切相连。因此，\BibleQuote{人若没有基督的圣神，就不是属祂的。}\newline{}\BibleRef{罗 8:9} 因此，唯独祂适当地荣耀主，因为经上说，\BibleQuote{祂要荣耀我，}\newline{}\BibleRef{若 16:14}不是作为受造物，而是作为\BibleQuote{真理之神}\newline{}\BibleRef{若 14:17}，在自己内清楚地显明真理，并且作为\Quote{智慧之神}，以祂自己的伟大启示\BibleQuote{基督是天主的德能和天主的智慧。}\newline{}\BibleRef{格前 1:24} 作为\OriginalTerm{护慰者}{Paraclete}，祂在自己内表达了差遣祂的那位护慰者的良善，并以自己的尊贵显露出祂所出自的那位的威严。因此，有一种本性的荣耀，如同光是太阳的荣耀；另一种是司法性的、自由意志\Emph{自外而来}的荣耀，施予那些配得的人。后者是双重的。经上说：\BibleQuote{儿子尊敬父亲，仆人敬畏主人。}\newline{}\BibleRef{拉 1:6} 这两者中，一种是奴仆的，由受造物给出；另一种可称为亲密的，由圣神完成。因为，正如主论及自己所说的：\BibleQuote{我在地上已经荣耀了你；你交托我的事，我已完成，}\newline{}\BibleRef{若 17:4} 同样，论及护慰者，祂说：\BibleQuote{祂要荣耀我，因为祂要从我那里领受，并要显给你们。}\newline{}\BibleRef{若 16:14} 并且，正如子被父荣耀，当父说\BibleQuote{我已经荣耀了\Emph{它}，还要再荣耀\Emph{它}，}\newline{}\BibleRef{若 12:28} 同样，圣神也通过祂与父和子的共融，以及独生子的见证而被荣耀，当祂说：\BibleQuote{人一切的罪和亵渎，都可得赦免；惟独亵渎圣神，总不得赦免。}\newline{}\BibleRef{玛 12:31}\par

47. 当我们藉着那光照我们的能力，凝视不可见天主的肖像之美，并通过这肖像被引向原型景象的至极之美时，我确信，知识的圣神就与我们不可分离，祂亲自将那瞻仰肖像的能力赐给爱慕真理景象的人，不是从外部展示，而是在祂内引领人达到圆满的知识。\Quote{除了圣子，没有人认识圣父。}因此，\Quote{除非受圣神感动，也没有人能说耶稣是主。}\BibleRef{格前12:3}因为经文不是说\Quote{藉着圣神}，而是\Quote{因圣神}；并且\Quote{天主是神，朝拜祂的人应当以心神以真理去朝拜祂。}\BibleRef{若4:24}正如经上所载\Quote{在你的光明中我们才能看到光明。}，即藉着圣神的光照，\Quote{那普照每人的真光。}\BibleRef{若1:9}结果是，祂在自己内显示独生子的荣耀，并在自己内将认识天主的知识赐予真正的朝拜者。因此，认识天主的途径是从一位圣神，通过一位圣子，到达一位圣父；反之，本性之善、内在之圣洁与王者尊严，从圣父通过独生者延伸到圣神。这样，既承认了\OriginalTerm{位格}{hypostases}，又不失关于\OriginalTerm{单一统权}{Monarchy}的真实教义。至于那些藉谈论第一、第二、第三来支持其\OriginalTerm{次位}{sub-numeration}的人，应当知道，他们是将异教的多种论引入基督徒纯洁的神学中。除了承认第一、第二、第三位天主之外，次位的可恶手段无法产生其他结果。对我们而言，主所规定的秩序已经足够。凡扰乱这秩序的人，其违法之罪绝不亚于不敬神的异教徒。\par

现在已经说了足够的话来证明，与他们的错误相反，本性的共融绝不会因次位的方式而被瓦解。不过，让我们向好争辩且软弱的对手让步，承认凡第二之物相对于某物而言被称为次位。现在让我们看看接下来是什么。\Quote{第一个人}，经上说，\Quote{出于地，属于土；第二个人是主，出自天上。}\BibleRef{格前15:47}又说：\Quote{先有的不是属神的，而是属生灵的；然后才是属神的。}\BibleRef{格前15:46}如果第二是相对于第一被列为次位，并且被列为次位者在尊严上低于它所列次位者，那么按照你们的说法，属神的在尊荣上低于属生灵的，天上的人低于地上的人。\par

% structure: p0112 source=h2 target=section reason=relative-heading-rank; base=h2

\section{第十九章}

\Emph{反对那些主张圣神不应受光荣的人。}\par

48. \Quote{就这样吧，}对方回答说，\Quote{但光荣绝不完全属于圣神，以致需要我们以赞美诗来高抬祂。}那么，如果祂与圣父及圣子的共融不被我们的对手视为有助于证明祂的地位，我们从哪里得到证明圣神尊严的证据，这尊严是\Quote{超越一切理解的}\BibleRef{斐4:7}呢？无论如何，我们有可能在一定程度上理智地领悟祂本性的崇高和祂不可接近的能力，通过观察祂称号的意义、祂作为的伟大，以及祂赐予我们——更确切地说，赐予所有受造物——的美好恩赐。祂被称为\Quote{圣神}，因为\Quote{天主是神}\BibleRef{若4:24}，以及\Quote{我们鼻孔的气息，主的受傅者。}祂被称为\Quote{圣}，\BibleRef{若一1:20}如同圣父是圣的，圣子是圣的；因为对于受造物，圣德是从外面引入的，但对于圣神，圣德是本性的圆满，因此祂被描述不是被圣化，而是圣化者。祂被称为\Quote{善}，如同圣父是善的，那从善者所生的是善的；对于圣神，祂的善就是本质。祂被称为\Quote{正直}，因为\Quote{主是正直的}，因为祂本身就是真理，本身就是正义，\BibleRef{格后3:8}没有偏差，也不偏斜，因为祂的实体不变。祂被称为\Quote{护慰者}，如同独生者一样，正如祂自己所说：\Quote{我要请求父，祂将赐给你们另一位护慰者。}因此，圣神与圣父、圣子共用这些名字，祂从祂本性的亲密关系中获得这些称号。还能从哪里得到它们呢？祂又被称为\Quote{君王}、\Quote{真理之神}、\Quote{智慧之神}。\BibleRef{依11:2}\Quote{天主的圣神}经上说\Quote{造了我}\BibleRef{约33:4}，天主又以\Quote{智慧、聪明和知识的神圣之灵}充满了比撒列。这样的名字是极其崇高和强大的，但它们并没有超越祂的光荣。\par

49. 祂的行动是什么呢？论尊威，不可言喻；论数目，不可胜数。我们怎能构想超越世代的事物呢？在我们可以想象的受造界存在以前，祂的行动是什么？祂赐予受造界何等大的恩宠？祂对于未来世代的权能如何？祂存在；祂先存；祂在万世之前与圣父和圣子共存。因此，即使你能想到任何超越世代的事物，你仍会发现圣神还在更高更远之上。若你想到受造界，天上的能力是由圣神建立，这建立指的是不能堕离善。因为正是来自圣神，能力才有与天主的亲密关系、不能转向邪恶以及恒久于福乐。是基督的来临吗？圣神是先驱。是降生成人的临在吗？圣神不可分离。奇迹的工作和治病的恩赐是藉着圣神。魔鬼因天主的圣神被驱逐。恶魔因圣神的临在而归于无有。罪的赦免是因圣神的恩赐，因为\Quote{你们被洗淨了，被圣化了，……因主耶稣基督之名，并因我们天主的圣神。}因圣神与天主有密切关系，因为\BibleQuote{天主派遣了祂圣子的圣神到你们心中，呼喊阿爸，父啊！}\BibleRef{迦 4:6} 死者复活藉着圣神的运作，因为\Quote{祢派遣祢的圣神，它们受造；祢又更新大地的面貌。}如果此处创造可理解为使离世者复生，那么圣神的运作何其伟大，祂为我们分赐复活后的生命，并使我们的灵魂契合于未来的属灵生命？或者如果此处创造指的是在今世已陷于罪中的人改变为更好的状态（因为圣经也如此使用，如\Person{保禄}的话：\BibleQuote{若有人在基督内，他便是新的受造物}\BibleRef{格后 5:17}），那在今世发生的更新，以及通过圣神在我们内从世俗感性的生活转变到天上的交往，那么我们的灵魂就被举扬到最高的钦佩境界。面对这些思想，我们是否害怕因过度尊崇而越过界限？难道我们不该惧怕，即使我们似乎给祂冠以人类思想所能构想或人言所能表达的最高名号，我们对祂的思念仍过于卑微吗？\par

是圣神说话，如同主说话一样：\BibleQuote{起来，下去和他们同去，不要疑惑，因为是我派他们来的。}\BibleRef{宗 10:20} 这是卑贱者的话，还是恐惧者的话？\BibleQuote{你们给我分出\Person{巴尔纳伯}和\Person{扫罗}，去做我召叫他们所做的工作。}\BibleRef{宗 13:2} 奴隶会这样说吗？\Person{依撒意亚}说：\Quote{主天主和祂的圣神派遣了我}，以及\Quote{圣神从上主那里降下，引导了他们。}但请不要再次把这种引导理解为某种卑微的服侍，因为圣言见证这是天主的工作——\Quote{祢领导了祢的子民}，经上说\Quote{如同羊群}，以及\Quote{祢领导了\Person{若瑟}如同羊群}，并且\Quote{祂安全地引导他们，使他们不惧怕。}因此，当你听见护慰者来了，祂要提醒你们，并\Quote{引导你们进入一切真理}时，不要曲解其意。\par

50. 但是，经上说：\BibleQuote{祂为我们转求。}\BibleRef{罗 8:26} 由此推论，既然恳求者低于恩主，那么圣神在尊严上也低于天主。但难道你没有听过关于独生子的话：\BibleQuote{祂在天主右边，也为我们转求}\BibleRef{罗 8:34} 吗？那么，不要因为圣神在你内——如果祂确实在你内——也不要因为祂教导本是瞎眼的我们，并引导我们选择有益的事，就不要因此容许自己被剥夺对祂正确而神圣的看法。因为以恩主的仁慈作为忘恩的借口，实在是不公平的极度妄为。\BibleQuote{不要使圣神忧郁，}\BibleRef{弗 4:30} 听听殉道者初果\Person{斯德望}的话，当他责备百姓的叛逆和悖逆时，他说：\BibleQuote{你们时常反抗圣神；}\BibleRef{宗 7:51} 再次，\Person{依撒意亚}说：\BibleQuote{他们使祂的圣神忧伤，因此祂转而为他们的仇敌；}\BibleRef{依 63:10} 在另一处：\Quote{雅各伯家激怒了主的圣神。}这些经文不是显明权能吗？我留待读者判断，当听到这些经文时，我们应当持何种观点：是把圣神视为工具、臣仆，与受造物同等，与我们同为仆人，还是相反，对虔诚者来说，这亵渎的低语难道不是极其可悲吗？你称圣神为仆人吗？但经上说：\BibleQuote{仆人不知道主人做的事，}\BibleRef{若 15:15} 然而圣神知道天主的事，正如\BibleQuote{人的心神，那在他里面。}\BibleRef{格前 2:11}\par

% structure: p0118 source=h2 target=section reason=relative-heading-rank; base=h2

\section{第二十章}

\Emph{反对那些主张圣神既非仆人亦非主人的等次，而是自由者等次的人。}\par

51. 据说，祂不是奴隶，也不是主人，而是自由的。哦，那些坚持此说者是何等可怕的麻木，何等可怜的胆大妄为！我是该为他们哀叹无知，还是该哀叹亵渎？他们试图将涉及天主性体的道理与人的道理相比较，从而侮辱前者，并企图将人间生活中等级差别可变的普通惯例应用于天主那不可言喻的性体，却未察觉在人间没有人生来就是奴隶。因为人要么因征服而受奴役的轭，如在战争中被俘；要么因贫穷而成为奴隶，如埃及人受\Person{法郎}压榨；要么，因一种明智而神秘的天命，最坏的子女被父亲命令罚去服事更明智和更好的人；任何正义的探究者都会宣称这不是定罪判决，而是一种益处。因为那因缺乏理智而自身没有自然统治原则的人，成为他人的财产，这更为有益，为使他受主人理性的引导，犹如战车有御者，船只舵手坐于舵柄。为此，\Person{雅各伯}因父亲的祝福成了\Person{厄撒乌}的主子\BibleRef{创 27:29}，为使那没有理智、即其正当守护者的愚昧儿子，即使不愿意，也能从他谨慎的兄弟那里获益。因此，\Person{客纳罕}将\BibleQuote{作他兄弟的仆人}\BibleRef{创 9:25}，因为他的父亲\Person{含}不智慧，他在德行上未受教导。那么，在现世，人就是这样成为奴隶，但那些逃脱了贫穷或战争，或不需要他人监护的人，是自由的。由此可见，即使一人被称为主人，另一人被称仆人，然而，鉴于我们彼此地位的平等以及作为造物主的财产，我们都是同等的奴隶。但在那另一个世界，你能从束缚中拿出什么？因为他们一受造，束缚就开始了。天体彼此不施行统治，因它们不为野心所动，但全向天主俯首，将作为主人应得的敬畏和作为造物主应得的荣耀同样归于祂。因为\BibleQuote{儿子孝敬父亲，仆人敬畏主人}\BibleRef{拉 1:6}，天主从所有人要求这两者之一；因为\BibleQuote{我若是父亲，我的孝敬在哪里？我若是主人，我的敬畏在哪里？}\BibleRef{拉 1:6}否则，所有人的生命若不在主人的监督之下，将是最可怜的；正如背叛的权势者的状况，他们因向全能的天主硬着颈项，挣脱了束缚的缰绳——并非他们的天性构造不同，而是由于他们对造物主的不服从性情。那么，你称谁为自由？那没有君王的人？那既无统治他人之权、也无被统治之意的人？在一切存在者中，找不到这样的本性。对圣神抱有这种概念，显然是亵渎。如果祂是一个受造物，当然祂与其余一切一同服事，因为经上说：\BibleQuote{万物都是你的仆人}；但如果祂超越受造界，那么祂就分享王权。\par

% structure: p0121 source=h2 target=section reason=relative-heading-rank; base=h2

\section{第二十一章}

\Emph{从圣经证明圣神被称为主。}\par

52. 但是，我们何必借争论这些不体面的问题来为我们的论点获取不公的胜利？我们本有能力藉更崇高的考量，证明荣光的卓越无可争议。如果我们真的重复\WorkTitle{圣经}教导我们的事，每一位\OriginalTerm{反圣神者}{Pneumatomachi}或许会高声喊叫，掩耳不听，捡起石头或任何随手可得的武器攻击我们。但我们不可将自身安危置于真理之前。我们从宗徒学会了：\Quote{主引导你们的心去爱天主，并耐心等待基督} \Emph{为我们的患难}。那引导你们去爱天主、为患难而耐心等待基督的主是谁？让那些想将圣神视为奴隶的人回答我们。因为如果论及天主圣父，它肯定会说：\Quote{主引导你们进入祂自己的爱}，或者如果论及圣子，它会加上\Quote{进入祂自己的忍耐}。那么让他们寻找还有哪一位相称受\Quote{主}的尊称。与此平行的是另一段经文：\Quote{并且愿主使你们彼此相爱的心，并爱众人的心，都能增长、充足，如同我们爱你们一样，好使你们的心在你们的天主父面前，在我们主耶稣基督同祂众圣徒来临时，成为圣洁，无可指责。} \BibleRef{得前 3:12} 现在，他恳求哪位主建立得撒洛尼信友的心，使他们在天主父面前，在我们主来临时，成为圣洁，无可指责？让那些将圣神置于奉差服役的仆役之列的人回答。他们不能。因此，让他们再听一个明确称圣神为主的见证。经上说：\Quote{主}就是\Quote{那圣神}；又说：\Quote{正如出于主的圣神}。但为了不留反驳的余地，我将引用宗徒的原话：\Quote{直到今日，在诵读\WorkTitle{旧约}时，同样的帕子还没有揭去；这帕子在基督内已被废去。……然而，当他们转向主时，帕子就被除去了。主就是那圣神。} 他为何这样说？因为那停留在文字表面的意义，并忙碌于法律仪文的人，他的心仿佛被犹太式的文字理解所包裹，如同帕子；他遭遇此事，是因为他不知道律法的肉身遵守已被基督的到来所废除，因为今后预象已转移到实体。太阳一出，灯就无用；真理显现，法律的职分消失，预言归于沉默。相反，那被赋予能力深入法律意义深处，经过文字的晦暗，如同经过帕子，进入不可言说之物的人，就像梅瑟与天主交谈时揭去帕子一样。他也从文字转向圣神。因此，梅瑟脸上的帕子对应于法律教导的晦暗，而属神默观对应于转向主。所以，在诵读法律时除去文字、转向主的人——主现在被称为圣神——更成为像梅瑟一样，他的脸因天主的显现而荣耀。因为正如靠近明亮颜色的物体自身也被周围的光泽所染，同样，那坚定定睛于圣神的人，藉圣神的光荣，仿佛变得更大辉煌，他的心好像被某种从圣神真理流出的光照亮。这就是\Quote{从圣神的光荣被改变成为祂自己的光荣}，不是吝啬的程度，也不是昏暗模糊，而是如同我们期望每一个被圣神光照的人那样。人啊，难道你不敬畏宗徒吗？他说：\Quote{你们是天主的宫殿，天主的圣神住在你们内。} \BibleRef{格前 3:16} 他岂能容忍用\Quote{宫殿}这样的称号来尊称奴隶的居所？那称\WorkTitle{圣经}为\Quote{天主所默感的}的人，因为\WorkTitle{圣经}是藉圣神的灵感而写成的，怎能使用侮辱和贬低祂的言语？ \BibleRef{弟后 3:16}\par

% structure: p0124 source=h2 target=section reason=relative-heading-rank; base=h2

\section{第二十二章}

\Emph{\Emph{从圣神与圣父、圣子同等地在思想上不可接近，建立其天然共融。}}\par

53. 此外，圣神性情的卓越超绝不仅可从祂与圣父、圣子共享同一称号并参与祂们的行动而得知，也可从祂如圣父、圣子一般，在思想上不可接近而得知。因为我们的主论及圣父超乎人的理解，论及圣子也是如此，祂也将同样的言语用于圣神。\Quote{正义的父}，祂说，\Quote{世界没有认识祢}，\BibleRef{若 17:25}此处以\Quote{世界}所指的，并非由天地构成的全体，而是我们这有死受无常的生命。当祂谈论自己时说：\Quote{还有不多的时候，世界不再看见我，你们却看见我}，\BibleRef{若 14:19}在此处，祂再次使用\Emph{世界}一词，指那些被物质和血肉生命束缚、仅凭肉眼认识真理的人，他们因不信复活，注定不再用心眼看见主。祂论及圣神也说了同样的话。\Quote{真理之神}，祂说，\Quote{世界不能接纳祂，因为看不见祂，也不认识祂；你们却认识祂，因为祂与你们同在}，\BibleRef{若 14:17}因为属血肉的人从未训练自己的心智默观，反而将它深埋在肉欲中如污泥一般，无力仰望真理的神性之光。因此，世界，即被肉性情欲所奴役的生命，不能接受圣神的恩宠，如同弱眼无法承受阳光一般。但是，主藉着祂的教导为纯洁的生活作证，将如今同时看见并默观圣神的能力赐给祂的门徒。因为\Quote{现在}，祂说，\Quote{你们因我给你们讲的话，已是洁净的了}，\BibleRef{若 15:3}因此\Quote{世界不能接纳祂，因为看不见祂，……你们却认识祂；因为祂与你们同在}，\BibleRef{若 14:17}依撒意亚也如此说：\Quote{铺张大地及其出产者，将气赐给其上之民，将神赐给践踏大地者}。因为那些践踏世俗事而超脱其上者，被见证为配得圣神的恩赐。那么，对于世界不能接纳、唯独圣人藉心灵的纯洁才能默观的祂，我们应当怎样看待？什么样的尊荣才可视为相称于祂？\par

% structure: p0127 source=h2 target=section reason=relative-heading-rank; base=h2

\section{第二十三章}

\Emph{列举祂属性的光荣化}。\par

54. 至于其他的权能，每一种都相信居于一个有限的地方。站在科尔乃略旁边的天使\BibleRef{宗 10:3}并未同时与斐理伯在一起；\BibleRef{宗 8:26}在祭台上与匝加利亚说话的天使，也并未同时占据他在天上的岗位。然而，圣神被认为同时运作在哈巴谷和在巴比伦的达尼尔身上，也曾与耶肋米亚同在地牢，与厄则克耳同在革巴尔河畔。\BibleRef{则 1:1}因为上主的神充满了世界，\BibleRef{智 1:7}且\Quote{我往何处能逃避你的神？我往何处能躲避你的面？}并且，按先知的话说：\Quote{因为我与你们同在，上主说……我的神仍在你们中间。}\BibleRef{盖 2:4}然而，对于那位无所不在、与天主同在的祂，我们应当赋予何种本性？是遍及万有的本性，还是局限于特定地方的本性，如同我们的论证所显示的天使的本性那样？没有人会这样说。那么，我们岂不该高举那位本性是神性、伟大无限、行动有力、恩赐良善的祂？我们岂不该给祂光荣？而我认为光荣不是别的，正是列举那属于祂的奇事。因此，要么我们的反对者禁止我们甚至提及从祂流到我们身上的美善，要么，仅仅重述祂的属性就是最圆满的光荣归与。因为即使是对于我们的主耶稣基督的天主和父，以及对于独生子，我们也不能以别的方式归光荣给他们，只能尽我们所能，重述属于他们的一切奇事。\par

% structure: p0130 source=h2 target=section reason=relative-heading-rank; base=h2

\section{第二十四章}

\Emph{从受造物中受光荣之物的比较，证明拒绝光荣圣神的荒谬}。\par

\Emph{此外}，人\Quote{以光荣和尊贵为冠冕}，并且\Quote{光荣、尊贵和平安}是应许\Quote{给一切行善的人}的。\BibleRef{罗 2:10}以色列人更有一种特殊的光荣，\Quote{他们承受了义子的名分、光荣……和礼仪}，\BibleRef{罗 9:4}圣咏作者也提到自己的光荣：\Quote{我的光荣要歌颂你}；又说\Quote{我的光荣，醒来吧！}按照宗徒，日月星辰各有光荣，并且\Quote{定罪的职分是有光荣的}。\BibleRef{格后 3:9}既然这么多事物都受了光荣，你难道希望只有圣神不受光荣吗？然而宗徒说\Quote{圣神的职分是有光荣的}。\BibleRef{格后 3:8}那么，祂自己怎会不配光荣呢？按照圣咏作者，义人的光荣是如何伟大，而按照你，圣神的光荣却等于零？难道这样的论证没有明显的危险，使我们陷入无可逃脱的罪吗？如果那因正义行为而得救的人尚且光荣敬畏上主的人，那么他岂不更不会剥夺圣神应得的光荣？\par

他们说，就算圣神应受光荣，但不可与圣父和圣子同受光荣。但有什么理由要放弃主为圣神所定的地位，而另创新说呢？有什么理由要从那位处处与天主性联合的圣神身上剥夺祂应得的荣耀呢？这荣耀体现在信仰宣认、救恩的圣洗、奇迹的施行、圣者心中的居留、以及赐予顺从者的恩宠上。因为没有圣神，就没有任何一项恩赐能达于受造物；甚至除非蒙圣神助佑，没有人能说出一句为基督辩护的话，正如我们从主及救主那里在福音中所学到的。\BibleRef{玛 10:19} 我不知道，任何领受过圣神的人，会同意我们忽视这一切，忘记祂在万物中的共融，并将圣神从圣父和圣子那里撕裂开来。那么，我们该把祂放在哪里，列于何等地位呢？与受造物同列吗？然而，所有受造物都处于束缚之中，而圣神却使人自由。\Quote{主的神在哪里，那里就有自由。}\BibleRef{格后 3:17} 我们可以提出许多论证，向那些认为将圣神与受造的本性并列是不相宜的人说明，但目前暂且略过。如果我要以与讨论的尊严相称的方式，提出我们方面常备的所有证据来推翻反对者的异议，则需要一篇冗长的论文，我的读者可能因我的繁冗而疲倦。因此，我打算将这个问题留待专题论文处理，而专注于我们眼前更紧迫的论点。\par

56. 那么，让我们逐一审视这些论点。祂本性是善的，正如圣父是善的，圣子是善的；而受造物则是通过择善而分享善。祂知道\BibleQuote{天主深奥的事；}\BibleRef{格前 2:10} 受造物藉着圣神领受不可言传之事的启示。祂与那产生并保存万有生命的天主，以及与那赐予生命的圣子一同赐予生命。经上说：\BibleQuote{那使基督从死者中复活的，也必要藉着那住在你们内的圣神，使你们有死的身体复活；}\BibleRef{罗 8:11} 又说：\BibleQuote{我的羊听我的声音，……我赐给他们永生；}\BibleRef{若 10:27} 但经上也说：\BibleQuote{圣神赐予生命，}\BibleRef{格后 3:6} 又说：\BibleQuote{圣神是生命，因为正义的缘故。}\BibleRef{罗 8:10} 主也证实说：\BibleQuote{使人生活的是圣神，肉一无用处。}\BibleRef{若 6:63} 那么，我们怎能将圣神与祂赋予生命的能力分离，使祂属于无生命的本性呢？谁如此好斗，谁如此完全没有天上的恩赐，未蒙天主美言的养育，谁如此在永生的希望中无分无份，以至于要将圣神从天主性中割裂，与受造物同列呢？\par

57. 现在有人极力主张，圣神在我们内是天主的恩赐，而恩赐不应受到与赐予者同等的尊敬。圣神是天主的恩赐，但却是生命的恩赐，因为经上说\BibleQuote{生命的圣神之律使我们获得了自由；}\BibleRef{罗 8:2} 也是能力的恩赐，因为\BibleQuote{当圣神降临于你们身上时，你们将领受能力。}\BibleRef{宗 1:8} 难道因此就应轻视祂吗？天主不也曾将祂的圣子作为无偿恩赐赐予人类吗？经上说：\BibleQuote{祂没有怜惜自己的儿子，反而为我们众人把祂交出来，岂不也把一切与祂一同赐给我们吗？}\BibleRef{罗 8:32} 另一处说：\BibleQuote{为使我们真正知道天主所白白赐给我们的一切，}\BibleRef{格前 2:12} 这是指降生成人的奥迹。由此可见，持这种论点的人，将天主仁爱的伟大作为亵渎的借口，实在比犹太人的忘恩负义有过之而无不及。他们挑剔圣神，因为祂赐予我们自由称天主为父。\BibleQuote{因为天主派遣了祂圣子的圣神到我们心中呼喊：\InnerQuote{阿爸，父啊！}}\BibleRef{迦 4:6} 使圣神的声音成为接受祂之人的声音。\par

% structure: p0136 source=h2 target=section reason=relative-heading-rank; base=h2

\section{第二十五章}

\Emph{圣经使用\Quote{在}或\Quote{藉着}一词（\OriginalTerm{在或藉着}{\GreekText{ἐν}}，见第3页注），代替\Quote{与}。其中也证明\Quote{和}一词与\Quote{偕同}具有相同效力。}\par

58. 然而，我们的对手却问：圣经为何从不说圣神与圣父和圣子一同受光荣，反而小心避免使用\Quote{与圣神}的表达，却处处更喜欢用\Quote{在祂内}来归光荣，认为这是更合适的说法？就我个人而言，我否认\Quote{在……内}（或\Quote{藉着}）一词比\Quote{与}一词暗示更低的尊威；相反，我认为，如果正确理解，它引导我们达到最高可能的意义。正如我们所观察到的，情形正是这样，它常常代替\OriginalTerm{与}{with}；例如，\Quote{我要以燔祭进入你的殿宇}，代替\OriginalTerm{与}{with}燔祭；\Quote{他也用金银带他们出来}，即\OriginalTerm{与}{with}金银；\Quote{你不在\OriginalTerm{在}{in}我们军队中出去}，代替\OriginalTerm{与}{with}我们军队；以及无数类似的经文。简而言之，我非常想从这种新奇的哲学中学习，按照我们对手所提出的从圣经而来的解释，宗徒用\Quote{在……内}一词归了什么荣耀，因为我从未找到\Quote{圣父啊，愿尊崇和荣耀藉着你的独生子，\OriginalTerm{藉着}{by}（或\OriginalTerm{在}{in}）圣神归给你}这样的格式——这种格式对我们对手来说，可以说是像呼吸空气一样自然。你确实可以分别找到这些分句，但他们绝无法向我们展示将它们如此联结在一起。因此，如果他们想要完全符合经文，就请他们给出准确的引用。另一方面，如果他们向习俗让步，就不能让我们成为这一特权的例外。\par

59. 既然我们在信友中发现两种表达方式都在使用，我们就两者都用；相信两种方式都同样给圣神以完全的荣耀。然而，最能堵住真理诋毁者之口的，是这个介词，它与圣经中的介词意义相同，但对我们对手来说不是那么得心应手的武器（事实上，它现在成了他们攻击的对象），并且它被用来代替连词\OriginalTerm{和}{and}。因为说\Quote{保禄和息耳瓦诺及弟茂德}\BibleRef{得前 1:1}与说保禄\OriginalTerm{与}{with}弟茂德和息耳瓦诺是完全相同的；因为两种表达方式都保持了名字的联系。主说\Quote{父、子和圣神}\BibleRef{玛 28:19}。如果我说父和子\OriginalTerm{与}{with}圣神，我是否在意义上作了任何改变？关于用连词\OriginalTerm{和}{and}连接名字的例子很多。我们读到\Quote{我们主耶稣基督的恩宠，和天主的爱，和圣神的相通}\BibleRef{格后 13:13}，以及\Quote{我因主耶稣基督，并因圣神的爱，恳求你们}\BibleRef{罗 15:30}。现在，如果我们想用\OriginalTerm{与}{with}代替\OriginalTerm{和}{and}，我们会作出什么区别？我看不出；除非有人根据严格的语法规则，更喜欢连词作为连接词并使之结合更强，而拒绝介词，认为它的力量较弱。但如果我们必须在这些点上为自己辩护，我想我们不需要长篇的辩护。事实上，他们的争论不是关于音节，也不是关于这个或那个词的声音，而是关于在能力和真理上极为不同的事物。正因如此，虽然音节的使用确实无关紧要，但我们的对手却努力授权某些音节，并将其他音节从教会中驱逐出去。就我个人而言，虽然这个词的功用一听到就很明显，但我仍将阐述导致我们的教父采取合理做法使用介词\Quote{\OriginalTerm{与}{with}}的理由。它确实与介词\Quote{\OriginalTerm{和}{and}}一样能驳斥萨贝里乌斯的祸害；它也能像\Quote{\OriginalTerm{和}{and}}一样清楚表明位格的区别，正如\Quote{我和我父要来}和\Quote{我与父原是一体}\BibleRef{若 10:30}这些话所显示的。此外，它所包含的永恒共融和不断结合的证明也是极好的。因为说圣子\OriginalTerm{与}{with}圣父同在，立即显明了位格的区别和共融的不可分性。同样的事甚至在纯粹人事中也可观察到，因为连词\Quote{\OriginalTerm{和}{and}}暗示行为中有共同元素，而介词\Quote{\OriginalTerm{与}{with}}在某种意义上也宣告了行为中的共融。例如——保禄和弟茂德航行到马其顿，但提希苛和敖尼息摩被派往哥罗森人那里。由此我们得知他们做了同样的事。但如果我们被告知他们\OriginalTerm{与}{with}某人航行，和\OriginalTerm{与}{with}某人被派遣？那么我们还得知他们是彼此一起进行该行动的。因此，虽然\Quote{\OriginalTerm{与}{with}}一词像其他任何词一样推翻了萨贝里乌斯的错误，它也击溃了犯完全相反方向错误的罪人；我是指那些以时间间隔将圣子与圣父分开，并将圣神与圣子分开的人。\par

60. 与\Quote{\Emph{在}}相比，存在这样的区别：\Quote{\Emph{同}}表示相关联者的相互结合——例如，那些同航、同住或以任何方式共同行事的人——而\Quote{\Emph{在}}则表示他们与其所从事之事的关系。因为我们一听到\Quote{\Emph{在}}或\Quote{\Emph{同}}这些词，便立刻想到船或房子。在日常用法中，这些词语的区别就是如此；费力的研究可能发现更多的例证。我没有时间检查音节的性质。既然已经表明\Quote{\Emph{同}}最清楚地表达了联合的意义，那么就让它被宣布为受保护之物，并停止对它发动野蛮而无休止的战争吧。尽管如此，虽然这个词本身是吉利的，但如果有人喜欢，在赞美颂中，用\Quote{和}这个音节来连接名字，并像我们在福音中所教导的，在圣洗的礼仪中给予光荣：\Quote{父、子及圣神}\BibleRef{玛 28:19}，那也随他：没有人会反对。在这些条件下，如果你愿意，让我们达成协议。但我们的敌人宁愿放弃他们的舌头，也不愿接受这个词。正是这一点激起了他们对我们的无情和无休止的战争。他们主张，我们应当\Emph{在}圣神内而非\Emph{和}圣神一起向天主献上荣耀的颂赞，并且他们热切地坚持\Emph{在}这个词，仿佛它贬低了圣神。因此，更详细地讨论它并非无益；我将感到惊讶，如果他们听到我们所坚持的，不把\Emph{在}本身当作他们事业的叛徒和圣神荣耀的逃兵。\par

% structure: p0141 source=h2 target=section reason=relative-heading-rank; base=h2

\section{第二十六章}

\Emph{\Quote{在}一词，无论有多少种含义，都是指圣神而言。}\par

61. 这表达虽然简短，但在我看来，深思之下，它的含义多种多样。因为我们发现，\Quote{\Emph{在}}一词的各种用法都有助于我们对圣神的理解。据说\Emph{形式}存在于\Emph{质料}之中；\Emph{能力}存在于能接受它的事物中；\Emph{习性}存在于受它影响的人中；等等。因此，既然圣神完美理性受造物，完成他们的卓越，祂类比于形式。因为那不再\Quote{按肉身生活}\BibleRef{罗 8:12}，而是\Quote{被天主的圣神引导}\BibleRef{罗 8:14}的人被称为天主的儿女，且\Quote{相似圣子的肖像}\BibleRef{罗 8:29}，被描述为属神的人。正如健康的眼睛有看的能力，同样，在净化了的灵魂中，有圣神的运作。因此保禄也为厄弗所人祈求，愿他们藉着\Quote{智慧的圣神}得以\Quote{眼目开通}\BibleRef{弗 1:17}。正如艺术存在于获得它的人身上，同样，圣神的恩宠也永远存在于领受者中，尽管不是持续运作。因为艺术潜在地存在于艺术家身上，但只有当他按艺术工作时才实际运作；同样，圣神总是与那些配得的人同在，但根据需要，在预言、医治或其他某种祂潜在行动的实际实施中运作。再者，如同健康、热或一般可变状态存在于我们身体内，同样，圣神也常常存在于灵魂内；因为祂不与那些因意志不稳而轻易拒绝所领受恩宠的人同住。例如在撒乌耳身上\BibleRef{撒上 16:14}，以及以色列子民的七十位长老（除了厄耳达得和默达得，只有他们似乎保留了圣神），以及任何与这些人性格相似的人。又如理性在灵魂中，有时是心中的思想，有时是口中发出的言语；同样，圣神也是如此，如祂\Quote{与我们的心神一同作证}\BibleRef{罗 8:16}，以及祂\Quote{在我们心中呼喊：阿爸，父！}\BibleRef{迦 6:4}，或者祂为我们说话，如经上说：\Quote{不是你们说话，而是你们父的圣神在你们内说话}\BibleRef{玛 10:20}。再者，在恩赐的分配上，圣神被理解为整体中的部分。因为我们彼此互为肢体，\Quote{按所受的恩宠各有不同的恩赐}\BibleRef{罗 12:5}。因此，\Quote{眼不能对手说：我不需要你；同样，头也不能对脚说：我不需要你}\BibleRef{格前 12:21}，而是所有肢体一起在圣神的合一中构成基督的身体，并互相提供由恩赐而来的必要帮助。\Quote{但天主随自己的意思将各肢体安置在身体上。}然而，\Quote{肢体彼此关怀}\BibleRef{格前 12:25}，按照同情的内在属灵共融。因此，\Quote{若一个肢体受苦，所有肢体一同受苦；若一个肢体受尊荣，所有肢体一同欢乐}\BibleRef{格前 12:26}。如同整体中的部分，我们个人也都在圣神内，因为我们众人\Quote{都受洗在一个身体内归于一个圣神。}\par

62. 这是一句非凡的陈述，但同样真实的是，圣神经常被称为成圣者的\Emph{地方}，而且将显而易见，即便通过这一比喻，圣神非但未被贬低，反而更受荣耀。因为为了清晰起见，圣经中常将适用于肉身的词语转用于属灵的概念。因此我们看到圣咏作者甚至在论及天主时说：\Quote{愿祢作我勇士的天主，作我的强垒以拯救我}，并论及圣神说：\Quote{看，在我旁边有一个地方，你要站在磐石上。}显然意指在圣神中的地方或默观，梅瑟进入那里后，得以领悟地看见天主向他显现。这是真正敬拜的特有而专属的地方；因为经上说：\Quote{你要谨慎，不可在任何地方献你的全燔祭……只可在上主你的天主所选择的地方。}\BibleRef{申 12:13}那么什么是属灵的全燔祭呢？\Quote{赞颂的祭献。}我们在什么地方献上它呢？在圣神内。我们从哪里学到这个？从主自己所说的话：\Quote{真正的朝拜者将用心神和真理朝拜父。}雅各伯看见这地方，说：\Quote{上主在这地方。}\BibleRef{创 28:16}由此可见，圣神实在是圣人的地方，而圣人是圣神适宜的地方，因为他将自己献作天主的居所，并被称为天主的圣殿。\BibleRef{格前 6:19}同样，保禄在基督内说话，说：\Quote{我们在天主面前，在基督内说话。}\BibleRef{格后 2:17}并且基督在保禄内，正如他亲自说的：\Quote{既然你们寻求基督在我内说话的凭证。}\BibleRef{格后 13:3}同样，他在圣神内讲论奥秘，\BibleRef{格前 14:2}而圣神也在他内说话。\BibleRef{伯前 1:11}\par

63. 因此，就受造者而言，圣神被说成\Emph{存在于}他们内，\Quote{多次并以多种方式，}\BibleRef{希 1:1}而就圣父和圣子而言，更符合真实敬拜的是主张祂不\Emph{存在于}，而是\Emph{与}他们同在。因为当祂居住在相称者内并执行祂自己的运作时，从祂流出的恩宠被恰当地描述为存在于那些能接受祂的人内。另一方面，祂在万世之前本质的存在，以及祂与圣子和圣父不间断的共处，若不使用表达永恒联合的称号便无法默观。因为绝对而真实的共存是针对那些彼此不可分离的事物而言的。例如，我们说热存在于热铁中，但在火本身的情况下，热是与火共存的；同样，健康存在于身体中，而生命与灵魂共存。由此可见，凡是亲密的、天生的、不可分离的共融，\Emph{与}这个词更具表达力，因为它暗示了不可分离的共融。另一方面，从圣神流出的恩宠自然地来去，则恰当而真实地说\Emph{存在于}，即使由于接受者坚定向善的习性，恩宠持续地与他们同在。因此，每当我们想到圣神本身的地位时，我们默观祂与圣父和圣子\Emph{同在}；但当我们想到从祂流出的恩宠在那些分享者内运作时，我们说圣神\Emph{存在于}我们内。而我们在圣神内所献的光荣颂，并非承认祂的地位；而是承认我们自己的软弱，同时表明我们无法靠自己光荣祂，我们的充足乃在于圣神。藉着祂被（在或藉）祂所推动，按照我们净化邪恶的程度，为所领受的恩惠向我们的天主献上感恩，因为我们每人领受圣神助佑的大小不同，好能献上\Quote{赞颂天主的祭献。}\BibleRef{希 13:15}据此一种用法，我们就这样按真实敬拜的要求在圣神内献上感恩；不过，若有人为自己作证说\Quote{天主的圣神在我内，我藉着从祂流出的恩宠获得智慧而献上光荣}，这并非完全没有异议。因为对保禄而言，说\Quote{我也想我拥有天主的圣神，}\BibleRef{格前 7:40}以及\Quote{要藉着住在我们内的圣神保管所托付给你的美好事物，}\BibleRef{弟后 1:14}是合适的。论及达尼尔，说\Quote{天主的圣神在他内}也是合适的，对于其他德行相似的人也是如此。\par

64. 但这短语也可另作解释：正如圣父在圣子内被看见，同样圣子也在圣神内被看见。所谓\Quote{在圣神内的敬拜}，暗示我们的理智是在光明中运作，这可以从主对撒玛黎雅妇人所说的话得知。她因本国习俗被误导，以为敬拜是地方性的；吾主为更好地教训她，说敬拜应当\Quote{以心神和真理}来举行\BibleRef{若 4:24}，其中\Quote{真理}明显地是指祂自己。因此，正如我们说敬拜是在天主圣父的肖像内，即是在圣子内，同样我们也说在圣神内的敬拜，因圣神在自己内显示了基督的天主性。因此，即使在我们的敬拜中，圣神也是与圣父和圣子不可分离的。你若停留在圣神之外，便根本无法敬拜；而当你进入祂内时，你绝不能将祂与天主分开——正如你不能将光明与可见物分开一样。因为若非藉着圣神的光照，便不可能看见不可见天主的肖像；而那注视肖像的人，也不可能将光明与肖像分开，因为视觉的原因必然与可见物同时被看见。因此，我们藉圣神的光照，适当地看见天主的\Quote{光荣的辉映}；又藉着\Quote{本体的真像}被引领到那真像和印鉴所印刻的祂那里。\par

% structure: p0147 source=h2 target=section reason=relative-heading-rank; base=h2

\section{第二十七章}

\Emph{论\Quote{偕}字的来源及其效力；并论教会的未成文法。}\par

65. 我们的反对者说，\Quote{\Emph{在}……内}这个词最贴切地适用于圣神，并且足以表达有关祂的一切思想。那么，他们问，我们为什么引入这个新短语，说\Quote{\Emph{偕}同圣神}而不说\Quote{\Emph{在}圣神内}，从而使用一个完全不必要的、且不为教会任何惯例所认可的表达呢？在本论文的前一部分，我们已经断言，\Quote{\Emph{在}……内}这个词并非特别分配给圣神，而是与圣父和圣子共通的。我也认为，已经充分证明，这个词非但没有减损圣神的尊严，反而将所有人——除非其思想完全扭曲——引向最高峰。现在，我要追溯\Quote{\Emph{偕}}字的来源，解释它有什么效力，并表明它与圣经是一致的。\par

66. 关于保存于教会中的无论是一般接受还是公开施行的信念与实践，我们有些来自成文教导，有些则是藉宗徒的传统以\Quote{奥秘}的方式传授给我们。这两者在关乎真宗教方面具有相同的力量。没有人会否认这一点——至少任何对教会的体制略知一二的人不会否认。如果我们试图以无权威为由拒绝那些没有成文规定的习俗，认为它们的重要性微不足道，我们就会无意中伤害福音的命脉；更确切地说，会使我们的公开定义沦为空洞的套语。例如，首先且最普遍的例子：谁曾以成文方式教导我们给那些信靠我主耶稣基督之名的人画十字圣号？哪篇作品教导我们在祈祷时转向东方？哪位圣人曾留下成文文字，记载了陈列圣体圣事面饼和祝福之杯时的祷词？因为我们并不满足于宗徒或福音所记载的，而是在前言和结论中都添加了其他对圣事有效性十分重要的言辞，这些来自不成文的教导。此外，我们祝福圣洗圣事的水和圣化圣油，并祝福受洗的慕道者。我们这样做有何成文根据？我们的根据岂非那无声的、神秘的传承吗？不仅如此，油的傅抹本身又是藉哪句成文教导的？三次浸入的习俗从何而来？至于其他洗礼习俗，我们从哪部圣经得到弃绝撒殚及其天使的教导？难道这不是来自父辈们以缄默守护的那未公开的、隐秘的教导，使其不受好奇干涉和刨根问底的研究吗？他们深知，奥秘的神圣尊严最好以缄默保存。连未入教者都不准看的东西，几乎不可能在成文文件中公开陈列。伟大的梅瑟为何不让所有人都看到会幕的各个部分？他将凡俗之人安置在圣所之外；将第一层庭院让给较洁净的人；他只认为肋未人配作天主的仆人；祭献、全燔祭及其他司祭的职务都分派给司祭；只从众人中选出一人进入至圣所，而且此人并非随时可以进去，只有一年中的一天，且那一天还定下了进入的时刻，使他得以凝视至圣所，为其奇特和新颖而惊叹。梅瑟有足够的智慧，知道轻蔑来自常见和显眼的事，而对不寻常和陌生的事则自然产生浓厚兴趣。同样，最初为教会立法的宗徒和教父们，以隐秘和缄默这样守护了奥秘的神圣尊严，因为随便向普通民众宣扬的事根本算不上奥秘。这就是我们保留不成文规诫和实践的传统的缘由，以免教理的知识因熟悉而被大众忽视和轻视。\OriginalTerm{信理}{Dogma}和\OriginalTerm{宣讲}{Kerugma}是两件不同的事；前者以缄默持守，后者向全世界宣讲。这种缄默的一种形式是圣经中使用的隐晦表达，使\OriginalTerm{信理}{dogma}的意义对读者难以理解，这对其极为有益：比如，我们都在祈祷时转向东方，但我们中很少有人知道我们是在寻找我们古老的家乡，即天主在东方伊甸园所栽植的乐园\BibleRef{创 2:8}。我们在每周的第一日站着祈祷，但并非所有人都知道原因。在复活（或\Quote{重新站起}；\OriginalTerm{复活}{\GreekText{ἀνάστασις}}）之日，我们通过站着祈祷提醒自己所获得的恩宠，不仅因为我们已与基督一同复活，并应当\BibleQuote{寻求上面的事}\BibleRef{哥 3:1}，而且因为这一天对我们而言，在某种意义上是我们期待之时代的预像。因此，虽然它是日的开始，梅瑟却不称它为\Emph{第一}日，而称它为\Emph{一}日。因为他说：\BibleQuote{过了晚上，过了早晨，是一日}，仿佛同一日经常重复。而\Quote{一}和\Quote{第八}是相同的，本身就是明确地指示那真正\Quote{一}和\Quote{第八}，即圣咏作者在某些圣咏标题中所提到的，那在现世之后的状态，那没有亏损或黄昏、后继者的日子，那没有终结或不老的时代。因此，教会必然教导其养子女在那日站着祈祷，以便通过持续提醒那无尽的生命，我们不会忽视为往那里迁移做准备。而且，整个五旬节是对未来世代所期待之复活的提醒。因为那\Quote{一}和\Quote{第一}日，如七乘以七，就构成了神圣五旬节的七个星期；因为从第一天开始，五旬节以同一天结束，通过类似的间隔日循环五十次。因此，它是永恒的相似，如环行路线，始于同一点并终于同一点。在这一天，教会的规则教导我们更偏爱直立的祈祷姿势，因为通过明确的提醒，它们仿佛使我们的心智不再停留于当下，而在于未来。此外，每当我们跪下并重新站起时，我们以行动表明，我们因罪恶而坠入地土，又因造物主的慈爱而被唤回天堂。\par

67. 倘若我试图追述教会那些未曾写下的奥迹，时间将不够我用。其余的我闭口不谈；但就连我们对圣父、圣子及圣神信仰的宣认，其文字来源是什么？如果承认，正如我们领受圣洗，同样也有责任去相信，我们按照我们圣洗的传统，并依照真正宗教的原则，以与我们圣洗相同的词语作出宣认，那么就让我们的对手也承认我们，在归光荣时像在宣认信仰时一样保持一致性。如果他们以缺乏文字权威为由反对我们的光荣颂，那就让他们给我们提供关于我们信仰宣认和我们列举的其他事项的文字证据。既然未写下的传统如此之多，且它们对\BibleQuote{虔诚的奥迹}\BibleRef{弟前 3:16}的意义如此重大，他们能拒绝我们使用一个从教父流传下来的词语吗？——我们发现在未败坏的教会中，它源于未经教导的习惯而存留着；——这是一个论据有力、对奥迹完整力量贡献不小的词语。\par

68. 现在已解释了两种表达的效力。我将再次说明它们在何处一致，何处彼此不同——并非它们相互对立，而是各自对真正宗教贡献其意义。介词\Quote{\Emph{在…内}}相对地更陈述关乎我们自身的真理；而\Quote{\Emph{与…一起}}则宣告圣神与天主的共融。因此我们使用两个词，藉前者表达圣神的尊威，藉后者宣告与我们同在的恩宠。这样，我们归光荣于天主，既\Quote{在}圣神内，也\Quote{与}圣神一起；在此我们并非使用自己的话，而是遵循主的教导，如同遵循固定的规则，并将祂的话转移到那些相连且密切关联、而在奥迹中的结合是必要的事物上。我们已认为自己有义务在我们的信仰宣认中，将那在圣洗中与祂们同列的一位结合起来，并将信仰宣认视为光荣颂的起源与父母。那么，该做什么？他们现在必须指导我们：要么不按所领受的方式施洗，要么不按所受洗的方式相信，要么不按所相信的方式归光荣。任凭任何人若能证明这些行为间的关系并非必要且不间断；或任凭任何人若能否认这里的创新必定意味着处处毁坏。但他们从不停止在我们耳边聒噪，说\Quote{\Emph{与}}圣神一起的光荣颂是未经授权、不合圣经等等。我们已经声明，就意义而言，说\Quote{愿光荣归于圣父、圣子\Emph{及}圣神}与说\Quote{愿光荣归于圣父、圣子\Emph{与}圣神一起}是相同的。任何人都不可能拒绝或取消\Quote{及}这个音节，它来自主亲口的话，也没有任何事物能阻止接受其等同者。两者之间有多少差异和相似，我们已展示过。我们的论证得到宗徒随意使用两个词的证实——他有时说\Quote{因主耶稣的名，并因我们天主的圣神}\BibleRef{格前 6:11}，有时说\Quote{当你们聚集时，连同我的心灵，并因我们主耶稣的大能}\BibleRef{格前 5:4}，并不认为使用连接词或介词会对名称的连接产生任何影响。\par

% structure: p0153 source=h2 target=section reason=relative-heading-rank; base=h2

\section{第二十八章}

\Emph{我们的对手拒绝将在圣经中用于人的术语，即与基督一同为王，用于圣神身上。}\par

69. 但让我们看看能否为教父的这种用法找到辩护；因为最先提出这种表达的人比我们自己更应受责备。保禄在致哥罗森人的书信中说：\BibleQuote{你们从前因罪恶和未受割损…祂已使你们一同}与基督\BibleRef{哥 2:13}。难道天主赐给整个民族和教会与基督同生的恩赐，而与基督同生却不属于圣神吗？但如果连想到这一点也是不虔敬的，那么按照本性紧密相连的天主三位如此宣认，岂不正是应有的敬畏吗？此外，这些人承认圣人们与基督同在（正如我们知道，保禄离开身体后就与主同在，去世后与基督同在），却尽其所能拒绝让圣神与基督同在，甚至不容许与人有同等的同在，这岂不显明他们缺乏无限的感觉吗？并且保禄在福音的职务中称自己为\BibleQuote{天主同工}\BibleRef{格前 3:9}；如果我们把\Quote{同工}这个词用于圣神——藉着祂，福音在天下一切受造物中结出果实——这些人会控告我们不虔敬吗？凡信赖主的人的生命\BibleQuote{已隐藏}，似乎\BibleQuote{与基督一同藏在天主内；当基督——我们的生命——显现时，}他们自己也要\BibleQuote{与祂一同在光荣中显现}\BibleRef{哥 3:3}。而生命之神本身，\BibleQuote{祂使我们脱离了罪的法律}\BibleRef{罗 8:2}，难道不与基督同在，无论是在与祂秘密隐藏的生命中，还是在我们将要显现在圣人身上的光荣的彰显中吗？我们是\BibleQuote{天主的承继者，基督的同承继者}\BibleRef{罗 8:17}，难道圣神在天主和祂的基督的共融中没有份吗？\BibleQuote{圣神亲自和我们的心神一同作证：我们是天主的子女；}难道我们不允许圣神享有我们从主所学的、祂与天主共融的见证吗？因为如果我们藉着在圣神内对基督的信德，希望与祂一同复活、一同坐在天上（当祂将我们卑贱的身体从自然变化到属灵时），却拒绝给圣神分享同坐、同荣以及我们从祂领受的一切，那就达到了愚昧的顶峰。在一切恩惠中，我们按照许诺者不可撤销的恩赐，相信自己堪当领受的，难道我们不让圣神分享任何一项，好像这一切都高于祂的尊严？按你的功德，你将\BibleQuote{永远与主同在}，你期望\BibleQuote{被提到云中，在空中与主相遇，这样我们就常与主同在}\BibleRef{得前 4:17}。你宣告那将圣神与圣父和圣子同列同序的人犯有不可容忍的不虔敬。你真的能否认圣神与基督同在吗？\par

70. 我羞于添加其余部分。你期望与基督同受光荣；（\BibleQuote{如果我们与祂一同受苦，也必与祂一同受光荣；}\BibleRef{罗 8:17}）但你却不与基督一同光荣\Quote{圣德之神}\BibleRef{罗 1:4}，好像祂不配与你同受尊荣。你希望\BibleQuote{与基督一同为王}\BibleRef{弟后 2:12}；但你却\BibleQuote{轻慢了恩宠的圣神}\BibleRef{希 10:29}，将祂置于奴隶和属下的地位。我说这话，不是要证明在光荣的归属中圣神应得多少，而是要证明那些连这一点也不愿给予、且像躲避不虔敬一样躲避圣神与圣子及圣父的共融之人的不公。谁能不叹息而谈论这些事？难道不是明明白白，连小孩子也能看出，目前的情况预示了信仰所受的威胁吗？无可否认的事变得不确定。我们宣认相信圣神，然后又与自己的宣认争吵。我们受洗，然后又开始争斗。我们呼求祂为生命之王，然后又像鄙视自己一样的奴隶鄙视祂。我们领受了祂与圣父和圣子同在，却又把祂当作受造物的一部分而羞辱。那些\BibleQuote{不知道该如何祈求}\BibleRef{罗 8:26}的人，即使被引导敬畏地说一句关于圣神的话，仿佛接近祂的尊严，却削减一切超出他们言语确切比例的内容。他们更应该哀叹自己的软弱，因为我们无力用言语表达对实际接受恩惠的感谢；因为祂\BibleQuote{超越一切理解}\BibleRef{斐 4:7}，并且指责言语天生无力丝毫接近祂的尊严；正如智慧书上所载：\Quote{你们要尽力高举祂，因为祂仍必超越；你们高举祂时，当用全力，不要疲倦，因为你们永远不能达及尽头。}的确，你们从不说谎的天主那里听到，对圣神的亵渎不得赦免\BibleRef{路 12:10}，对于这类言语，你们要交账的，这是可怕的。\par

% structure: p0157 source=h2 target=section reason=relative-heading-rank; base=h2

\section{第二十九章}

\Emph{列举教会中曾在其著作中使用\Quote{同}字的杰出人物。}\par

71. 针对\Quote{偕同圣神}形式的《光荣颂》没有成文权威的反对意见，我们坚持认为：如果任何不成文的事例都不可接受，那么此形式便不应被接纳。但既然我们习俗中的大部分奥秘都是没有成文权威而获接纳的，那么，就让我们与众多其他事例一同接纳此形式。因为我认为持守口传也是宗徒的作风。经上说：\BibleQuote{我称赞你们}，经上说，\BibleQuote{因你们在一切事上记念我，并持守我所传授给你们的规条}；\BibleRef{格前 11:2} 又说：\BibleQuote{你们要持守你们所学习的传统，无论是藉着言语，或是藉着我们的书信。}\BibleRef{得后 2:15} 这些传统之一便是我们眼前的做法，它由起初的制定者所建立，稳固地植根于各教会，传给他们的继承者，且经由长期使用，它与时间同步推进。倘若我们像在法庭上那样缺乏书面证据，却能在你们面前带来大量证人，你们难道不会投票判我们无罪吗？我想会的，因为\BibleQuote{凭两三个人的口作证，即可定案。}\BibleRef{申 19:15} 若我们能清楚证明长期的时间对我们有利，我们岂不显得有理，认为这诉讼不应在你们面前提出吗？因为古老的教义带有某种敬畏之感，因其白发苍苍的古老而可敬。因此，我将列出此词支持者的名单（也必须考虑时间因素，以说明它所历经的不容置疑的时期）。因为此词并非源于我们。如何能呢？与这个词流行的时间相比，我们，用约伯的话说，是\BibleQuote{昨日才有。}\BibleRef{约 8:9} 我自己——若必须谈及与我个人相关的事——珍视此短语，视之为我父辈留给我的遗产。它是由一位在事奉天主上度过漫长一生的人传给我的，我由他受洗，并被接纳进入教会的职务。在我尽我所能考察古时的蒙福之人是否使用现在受到反对的词语时，我发现了许多值得信赖的人，不仅因其年代久远，还因其一个与今日之人不同的特征——他们知识的精确性。他们中有些人在《光荣颂》中用介词连接此词，另一些人用连词，但从未被认为行事分歧——至少就真宗教的正确意义而言是如此。\par

72. 著名者如\Place{罗马}的\Person{依勒内}和\Person{克肋孟}；\Place{罗马}的\Person{狄约尼削}，以及——说来奇怪——\Place{亚历山大的}\Person{狄约尼削}，他在致同名者的第二封书信\WorkTitle{论确信与辩护}中如此结语。我将引证他的原话：\Quote{跟随所有这些圣贤，我们也从我们之前的司铎那里领受了一种形式和规则，与他们以同样的措辞感恩，以此结束致你们的书信。归光荣与权能于天主父及圣子我们的主耶稣基督，偕同圣神，至于无穷之世。亚孟。}无人能说这段文字被改动过。若他说的是\Quote{\Emph{在}圣神内}，他不会如此坚持说他领受了形式和规则。因为这一短语的使用很多，需要辩护的乃是\Quote{\Emph{偕同}}的使用。此外，狄约尼削在他的论文中间部分，为反对撒伯里乌异端者写道：\Quote{若因位格是三个而说它们被分开，那它们就是三个，尽管他们不喜欢。否则，就让他们彻底摧毁天主圣三。}又说：\Quote{因此最神圣的是在统一之后的天主圣三。}克肋孟以更原始的方式写道：\Quote{天主活着，主耶稣基督活着，圣神活着。}现在让我们听听依勒内——他生活在宗徒时代附近——在他的著作\WorkTitle{驳斥异端}中如何提到圣神：\BibleQuote{宗徒正确地称那些放纵无度、随从私欲、对圣神毫无渴望的人为\Emph{属血肉的}。}在另一处，依勒内说：\BibleQuote{宗徒大声宣告，血肉不能承受天国，免得我们因没有分沾神圣之神而达不到天国。}若有人因\Place{巴勒斯坦}的\Person{欧瑟伯}的丰富经验而认为他可信，我进一步指出他在讨论古人多妻问题时所使用的原话。他激励自己投入著作，写道：\Quote{呼求先知们的圣洁天主，光明之源，藉着我们的救主耶稣基督，偕同圣神。}\par

73. \Person{奥利振}也在不少圣咏释义中使用\Quote{\Emph{偕同}圣神}的赞美辞形式。他对圣神的看法并非始终如一，亦非处处正确；然而在许多段落中，甚至连他自己也虔敬地承认既定习惯的力量，并以符合真宗教的言辞论及圣神。如果我没有弄错，正是在他的\WorkTitle{若望福音释义}第六卷中，他明确地将圣神当作敬拜的对象。他的话是这样的：\Quote{洗涤或水是灵魂洁淨的象征，灵魂被洗去一切因邪恶而来的污秽；但就它本身而言，对那将自己交付于可钦崇的天主圣三之神性的人，藉着呼求的能力，它也是恩惠的根源与泉源。}此外，在他的\WorkTitle{罗马书释义}中，他说\Quote{神圣的权能，}\Quote{能够接纳独生者，以及圣神的神性。}因此，我认为，传统的有力影响常常迫使人们表达出与他们自己观点相矛盾的话语。此外，这种赞美辞的形式甚至对历史学家\Person{阿夫里加努斯}也并非陌生。在他的\WorkTitle{时代概要}第五卷中，他说：\Quote{我们深知这些术语的分量，且未忽略信德的恩宠，我们感谢圣父，祂将祂自己的受造物、世界的救主和我们的主耶稣基督赐予我们，愿光荣与尊威偕同圣神归于祂，至于永远。}其余的段落或许令人怀疑，或者可能确实已被改动，而且被篡改的事实难以察觉，因为差异仅在于一个音节。然而，我所全文引用的那些段落，则无法被任何不诚实的手脚所触及，并且容易从实际著作中得到验证。\par

我现在要提出另一项证据，它或许看似微不足道，但因其古老，绝不能被一个被指控标新立异的被告所忽略。我们的祖先认为，在黄昏时分不应沉默地接受光明的恩赐，而应在光明显现时立即感谢。这些点灯时的感谢之词出自何人之手，我们无法断言。然而，民众仍念诵着这古老的格式，从未有人将说\Quote{我们赞美圣父、圣子，以及天主的圣神}的人视为不敬。若有人知道\Person{阿特诺格内斯}的\WorkTitle{赞歌}——他在匆匆走向烈火成全时，将其作为某种告别礼物留给朋友——他就知道殉道者对圣神的看法。关于这一点，我不再多说。\par

74. 然而，我将把伟大的\Person{额我略}和他的言论置于何处？难道我们不应将这样一个人置于宗徒和先知之列吗？他曾与他们同一位圣神同行\BibleRef{格后 12:18}，一生从未偏离圣人的足迹，在世时始终持守福音公民的精确原则。我确信，如果我们拒绝将那颗灵魂算作天主的子民，就是对不起真理；它在天主的教会中如灯塔般闪耀。因为藉着圣神的协同运作，他制伏魔鬼的力量是惊人的；他蒙受神恩，拥有\BibleQuote{为服从信德，在万民中}的言语恩宠\BibleRef{罗 1:5}，以至于尽管只有十七位基督徒被交付给他，他却使所有城乡民众藉着认识归向天主。他亦藉基督大能的名号命令河流改道，使一个曾引起一些贪婪弟兄争吵的湖泊干涸。此外，他对未来之事的预言毫不逊色于伟大的先知们。详细叙述他所有的奇事将过于冗长。由于圣神以一切权能、奇迹和异事在他内所行的丰盛恩赐，他甚至被教会的外敌称为第二位梅瑟。因此，凡他所藉恩宠成就的一切，不论言语或行为，似乎总有一道光芒照耀，那来自不可见之天上的权能的标志常伴随着他。直到今日，他仍是本乡人民极其敬仰的对象；他在教会中建立的记忆永远鲜活翠绿，不为时间所磨灭。因此，除了他所遗留给教会的，没有任何习惯、言论或奥迹礼规被添加到教会中。所以，确实因其制度的古老，他们的许多礼仪似乎有欠缺。因为他在管理教会上的继承者们无法容忍在其所认可之外接受任何后来的发现。如今，额我略的制度之一正是如今遭到反对的赞美辞形式，教会凭其传统的权威得以保存；任何愿意短途旅行的人都不难验证这一说法。我们的\Person{斐尔米利亚诺}持有这一信念，他所留下的著作可以作证。著名的\Person{默肋提乌斯}的同代人亦说他怀有这种意见。但为何引用古代的权威？如今在东方，真宗教的维护者们难道不正是凭这一措辞而被众所周知，并像使用口令一般与他们的对手区分开吗？我曾听一位\Person{美索不达米亚人}——此人精通当地语言且观点纯正——说，按他本国的习惯，任何人即便愿意也不可能以其他方式表达，他们被母语的惯用法迫使用\Quote{与}这个音节来献上赞美辞，或者更准确地说，用等同的表达。我们\Person{卡帕多西亚人}也用本地方言这样说，圣神早在语言分开时就预见了这一短语的用处。整个西方呢？从\Place{伊利里亚}几乎到我们世界的边界，不也支持这一措辞吗？\par

75. 那么，我怎会是新名词的革新者和创始人呢？我提出整个民族、城市、远超人类记忆的古老习俗、以及教会柱石和以全备知识与属神能力著称之人作为这词语的起源与护卫者。因此，这结盟的敌军向我发动攻击，城镇、村庄和最遥远的地区都充满了毁谤我的人。这些事对寻求和平者是悲伤和痛苦的，但为信德而受的苦难将得到极大的忍耐赏报。此外，让刀剑闪耀，让斧头磨利，让火焰比巴比伦的更烈，让一切刑具都向我发动。对我来说，没有什么比不惧怕主针对亵渎圣神者所发出的威胁更可怕的了。善意的读者会在我所说的话中找到满意的辩护：我接受一个如此为圣人们所珍爱和熟悉的、并经如此长久使用所确认的词语，因为从福音首次宣讲直到我们时代，它已被证明在教会内享有完全的权利，并且最重要的是，它被接受为带有符合圣洁与真信仰的意义。但在大审判台前，我准备了什么作为我的辩护呢？这就是：我首先被主在圣洗圣事中使自己与祂和祂的父联合所赐予的尊荣引向圣神的光荣\BibleRef{玛 28:19}；其次，通过这种入门礼，我们每个人都被引入对天主的认识；最重要的是，对受威胁惩罚的恐惧排除了所有不敬和不相称的思想。但我们的对手会说些什么呢？他们既不敬畏主的尊荣，也不畏惧祂的威胁，那么他们将为他们的亵渎作何辩护？他们需要决定自己的行为，或者即使现在也要改变。就我而言，我最恳切地祈求善心天主使祂的平安在我们所有人心中作主，使这些骄傲膨胀、向我们列阵的人因温和与爱德之神而平静下来；如果他们变得完全野蛮、不可驯服，愿祂至少赐予我们恒久忍耐，忍受我们不得不从他们手中承受的一切。总之，\Quote{那在自己心里有死刑判决的人}\BibleRef{格后 1:9}，为信德而受苦并不痛苦；难以忍受的是不能为信德而战。运动员在比赛中受伤的抱怨，不如连入场资格都争取不到的痛苦。或者，这也许是智慧者撒罗满所说的沉默之时\BibleRef{训 3:7}。因为，当生命被如此猛烈的风暴冲击，以至于所有受教于圣言的人的心智都被错误推理的欺骗充满、如同眼睛被尘土迷住时，当人们被奇特而可怕的声音震晕，当整个世界摇动、一切摇摇欲坠时，像我这样对着风呼喊又有什么益处呢？\par

% structure: p0165 source=h2 target=section reason=relative-heading-rank; base=h2

\section{第三十章}

\Emph{论教会目前状况}\par

76. 那么，我将我们目前的情况比作什么呢？我想，它可以比作一场因旧日争执而起的海战，交战双方都怀有刻骨仇恨，久经海战，求战心切。看，我恳请你，在你眼前呈现这幅画面。看那敌对舰队以可怖的阵列冲向攻击。它们以一股无法控制的狂怒交战，决一胜负。你大可以想象，船只被狂风暴雨驱赶得来回飘荡，浓云密雾使整个场景一片漆黑，口令在混乱中无法分辨，敌友难分。为了充实这幅想象画面，设想大海波涛汹涌，从深渊翻腾，同时暴雨如注从天而降，可怕的巨浪高涌。从天上四面八方刮来的风都集中打在一处，两军舰队互相撞击。一些战士背叛，一些在激战中逃亡，一些必须同时迎风催舟和抵御袭击者。对权威的嫉妒和对个人主宰的贪欲使水手们分裂成互相残杀的各派。除此之外，还要想想那弥漫整个海上的混乱而无意义的喧嚣——呼啸的风、破碎的船、沸腾的浪、交战者用各种声音表达不同情感的喊叫，使舰队司令或舵手听不到一句话。混乱和骚动巨大无比，因为当生命绝望时，极端的苦难给了人任意作恶的自由。此外，设想所有的人都患上了贪求光荣的疯狂不治之症，即使自己的船正沉入深渊，他们也不停止互相争胜。\par

77. 请现在，我恳求你，从这象征性的描述转向那令人不快的现实。难道不是一度出现，亚略异端在分裂成一个与天主的教会对立的派别后，独自以敌对阵势站立？但是，当我们的敌人对我们的态度从长期苦战转为公开战争时，众所周知，战争分裂成了我无法尽述的许多分支，以致所有人都被共同的党派精神和个人的猜疑激起了根深蒂固的仇恨。有哪一场海上的风暴曾像这场教会的暴风雨一样凶猛狂暴？在其中，教父们的一切界标都被挪移；每一种根基、每一道意见的壁垒都被动摇；一切浮在不可靠之物上的东西都被摔碎和摇落。我们互相攻击，我们互相推翻。如果敌人没有先击打我们，我们就被旁边的同伙所伤。如果敌人被击倒，他的战友便将他踩在脚下。我们之间至少有一个联合的纽带，就是我们仇恨共同的敌人，但敌人一过去，我们便在彼此中找到了敌人。谁能完全列举所有的残骸？有的在敌人攻击下沉没，有的因盟友的意外背叛，有的因自己的军官失策。我们看到，仿佛整个教会，连同船员全部，撞碎在虚伪异端的暗礁上，而另一些敌视\OriginalTerm{救恩之神}{Spirit of Salvation}的人却夺取了船舵，使信仰搁浅。\BibleRef{弟前 1:19} 然后，这世界的有权者所造成的扰乱\BibleRef{格前 2:6}，以飓风或旋风都无法比拟的暴力，导致了民众的堕落。天主设立来光照人们灵魂的世界光体，被赶出了自己的家园，一种真正阴郁沮丧的黑暗笼罩了教会。普遍毁灭的恐怖已然逼近，但他们彼此间的竞争却如此无度，以致钝化了所有危险感。个人的仇恨比普遍共同的战争更为重要，因为那些认为即刻满足野心比末世等待着我们的赏报更宝贵的人，宁愿选择胜过对手的荣耀，而不顾人类共同的福祉。所以，所有人都一样，各尽其能，互相举起杀人的手。争论中相互交锋的斗士发出刺耳的喊声；整个教会几乎充满了无意义的尖叫、不可理解的噪音，这些噪音来自不停地搅动，使真实宗教教义的正确规则转向，一时过度，一时不足。一方面，有些人混淆了位格，被带入犹太教；另一方面，有些人因本性的对立而堕入异教。在这两派之间，受默感的圣经无力调解；宗徒的传统无法提出仲裁条款。直言不讳有害友谊，意见分歧就是争吵所需的一切理由。没有一种联盟的誓言能像错误的相似那样有效地使人忠于叛乱。每个人都是神学家，尽管他的灵魂上有数不清的斑点。结果是，创新者发现大量适合派别争斗的人，而自封的猎位者家族的子孙拒绝圣神的治理，瓜分了教会的主要尊位。福音的体制如今因缺乏纪律而到处混乱；各人争先恐后地抢夺高位，每个自我宣传者都试图强行进入高职。这种命令欲的结果是，我们的子民因缺乏秩序而陷入狂乱状态；掌权者的劝诫变得完全无目的和无效，因为没有一个人不因愚昧的厚颜无耻而认为，发号施令既是自己的职责，也是服从别人的职责。\par

78. 所以，既然在这样的混乱中没有任何人声足够响亮，我认为沉默比言语更有益，因为如果训道者的话有任何真理：\Quote{智慧人的言语在安静中被听见}\BibleRef{训 9:17}，那么在目前的情况下，谈论它们必定一点也不合宜。此外，先知的话也约束我：\Quote{因此，智慧人在那时必静默不言，因为这是邪恶的时代}\BibleRef{亚 5:13}。这个时代，有人绊倒邻舍的脚跟，有人践踏倒下的人，还有人欢乐拍手，却没有一人同情跌倒者并伸出援手，尽管按照古法，连看见仇人的驴在重担下跌倒而经过的人，也不得无罪。\BibleRef{则 23:5} 但现今并非如此。为何不是？许多人的爱心冷淡了；\BibleRef{玛 24:12} 兄弟间的和睦被摧毁，合一之名被无视，兄弟的劝诫不再被听见，基督的怜悯无处可寻，同情的眼泪不再流下。如今没有一人接待\Quote{信德软弱的}\BibleRef{罗 14:1}，但同族人之间的相互仇恨如此高涨，以致他们更因邻人的跌倒而欢喜，胜过自己的成功。正如在瘟疫中，生活最规规规矩矩的人也染上同样的疾病，因为他们通过与感染者接触而传染，如今因这占据我们灵魂的邪恶竞争，我们被卷入效法邪恶中，我们全都彼此一样坏。因此，对犯错者的审判是无情而苛刻的；对善良者的批评是无情而敌对的。这邪恶在我们中间扎得如此之深，以致我们变得比野兽更野蛮；野兽尚且与同类群居，而我们最残酷的战争却是对自己人。\par

79. 由于这一切理由，我原该保持沉默，但爱却将我拉向另一方向，它\Quote{不求己益，}\BibleRef{格前 13:5}并渴望克服时间和环境所设置的一切困难。我也从巴比伦的孩童那里学到：当无人支持真宗教的事业时，我们应当单独且无援地完成我们的责任。他们从火焰中高声赞美歌颂天主，不顾那轻视真理的群众，虽仅三人，却彼此同心。因此，我们也不为敌人的乌云所惧，将希望寄托于圣神的助佑，勇敢地宣讲了真理。倘若我没有这样做，那真是可怕的事：亵渎圣神的人竟如此轻易地变得胆大妄为，攻击真宗教；而我们，拥有如此强大的盟友和支持者，却退缩，不敢为那由教父传统、藉着不断的记忆传递至今的教义服务。推动我进行这项工作的另一强大动力，是你们\Quote{无伪的爱，}的炽热，以及你们性情中的严肃与静默；这保证你们不会将我所说的话公之于众——不是因为不值得宣扬，而是为了避免把珍珠丢在猪前。\BibleRef{玛 7:6}我的任务现已完成。如果你们对我的阐述感到满意，就让我们就此结束讨论。如果你们认为某一点需要进一步阐明，请毫不犹豫地勤加探究，并通过提出任何无争议的问题来增加你们的认识。主必会通过我或他人，就尚待阐明的事赐予圆满的解释，按圣神赐予配得者的知识。阿们。\par

\SourceRecord{Translated by Blomfield Jackson.；Nicene and Post-Nicene Fathers, Second Series；Vol. 8.；Edited by Philip Schaff and Henry Wace.；Buffalo, NY: Christian Literature Publishing Co.,；1895.；<http://www.newadvent.org/fathers/3203.htm>.}
